% ---- exact/opening ----
\documentclass[11pt]{amsart}
\usepackage{amsmath,amssymb,amsthm,mathtools,booktabs}
\usepackage[margin=1in]{geometry}
\usepackage{xurl}
\usepackage[colorlinks=true,linkcolor=blue,citecolor=blue,urlcolor=blue]{hyperref}
\setlength{\emergencystretch}{3em}
\newtheorem{theorem}{Theorem}[section]
\newtheorem{proposition}[theorem]{Proposition}
\newtheorem{lemma}[theorem]{Lemma}
\newtheorem{corollary}[theorem]{Corollary}
\newtheorem{maintheorem}{Theorem}
\renewcommand{\themaintheorem}{\Alph{maintheorem}}
\theoremstyle{definition}
\newtheorem{definition}[theorem]{Definition}
\theoremstyle{remark}
\newtheorem{remark}[theorem]{Remark}
\newcommand{\N}{\mathbb N}
\newcommand{\Z}{\mathbb Z}
\newcommand{\Q}{\mathbb Q}
\newcommand{\R}{\mathbb R}
\newcommand{\E}{\mathbb E}
\newcommand{\PP}{\mathbb P}
\newcommand{\Ztwo}{\mathbb Z_2}
\newcommand{\Qtwo}{\mathbb Q_2}
\newcommand{\haar}{h}
\newcommand{\one}{\mathbf 1}
\DeclareMathOperator{\dist}{dist}
\DeclareMathOperator{\supp}{supp}
\DeclareMathOperator{\Var}{Var}
\title[Exact rates in Collatz synchronization]{Exact rates in Collatz synchronization}
\author{David Leen}
\email{daleen@gmail.com}
\date{October 2026}
\subjclass[2020]{11B83, 37A44, 60G50, 11K38}
\keywords{Collatz map, synchronization, Gao's conjecture, absorption time,
two-adic Haar measure, occupation measure}
\begin{document}
\begin{abstract}
Gao conjectured that his dyadic proportions of coalescing adjacent
integers tend to one. We determine the exact deficiency:
\[
1-\bar d_k\sim\frac{\beta}{\sqrt k},\qquad
\beta=\frac{\sqrt6\,\E M}{\sqrt\pi},\qquad 1\le\E M<\infty.
\]
Here $M$ counts the excursions of the canonical adjacent-pair coupling.
The balanced shortcut absorption time has coefficient $2\E M/\sqrt\pi$;
the first ordinary meeting time has coefficient $\beta$. The same
$\beta$ governs the fraction of integers $n<N$ with no ordinary meeting
by $\lfloor\log_2n\rfloor$.
The proof passes from stationary rigidity to actual killed occupation
and then determines the physical clock of a long excursion. A weighted
chronological assembly lemma handles dependent excursions using only
their expected total weight. This weight also yields the harmonic
survival capacity from every integer canonical state and extends the
coefficient to the exact admissible offset domain $\Z[1/3]$.
A uniform finite-horizon estimate gives synchronized growing-block
versions of Korec's and Tao's descent theorems. No termination
assumption is used.
\end{abstract}
\maketitle

\section{The rates and their common mechanism}\label{sec:intro}

For the ordinary and shortcut maps write
\[
C(n)=\begin{cases}n/2,&n\text{ even},\\3n+1,&n\text{ odd},\end{cases}
\qquad
T(n)=\begin{cases}n/2,&n\text{ even},\\(3n+1)/2,&n\text{ odd}.\end{cases}
\]
They also act on $\Ztwo$, with parity modulo two. Set
$\N=\{0,1,\ldots\}$; unmarked logarithms are natural. Define
\[
o_t(x)=\#\{0\le j<t:T^j(x)\text{ odd}\},\qquad
H(x,m)=\#\{0\le j<m:C^j(x)\text{ even}\}.
\]
Under a Haar root $x\in\Ztwo$, the adjacent clocks are
\begin{align*}
A&=\inf\{t:T^t(x)=T^t(x+1),\ o_t(x)=o_t(x+1)\},\\
B&=\inf\{m:C^m(x)=C^m(x+1),\ H(x,m)=H(x+1,m)\},\\
U&=\inf\{m:C^m(x)=C^m(x+1)\}.
\end{align*}
Empty hitting sets have infimum $\infty$. Equality of both endpoint and
count is called balanced merger. The canonical coupling orients its
entries as $Y=3^sX+D$, $s\ge0$, and absorbs at $(s,D)=(0,0)$.
Let $M$ count \emph{all} odd-offset level-zero departures before
absorption, including departures inside finite absorbing-word trials.
The initial state $(0,1)$ departs, so $M\ge1$.

For positive integers, write $G(n,k)$ if some $m\le k$ satisfies
\[
C^m(n)=C^m(n+1),\quad H(n,m)=H(n+1,m),\quad
C^j(n),C^j(n+1)>1\quad(0\le j\le m).
\]
Gao's original proportion is
\[
\bar d_k=2^{-k}\#\{1\le n<2^k-1:G(n,k)\}.
\]
The strict interval has $2^k-2$ starts but denominator $2^k$.
The different logarithmic-deadline population is
\[
\mathcal G=\{n\ge1:\exists j\le\lfloor\log_2n\rfloor,
 C^j(n)=C^j(n+1)\},\qquad
f_N=1-\frac{\#(\mathcal G\cap[0,N))}{N}.
\]

\begin{maintheorem}[Exact synchronization rates]
\label{thm:exact}\label{thm:gao-sharp}\label{thm:deadline}
Put
\[
\kappa=\frac{2\E M}{\sqrt\pi},\qquad
\beta=\sqrt{\frac32}\,\kappa=\frac{\sqrt6\,\E M}{\sqrt\pi},
\qquad 1\le\E M<\infty.
\]
Then
\[
\sqrt t\,\PP(A>t)\to\kappa,\qquad
\sqrt t\,\PP(B>t)\to\beta,\qquad
\sqrt t\,\PP(U>t)\to\beta,
\]
and, for Gao's original predicate and all integer cutoffs $N$,
\[
\sqrt k(1-\bar d_k)\to\beta,\qquad
\sqrt{\log_2N}\,f_N\to\beta.
\]
On finite balanced absorption, $U=B=A+o_A(x)-1$ and both paths have
$A-1$ halvings at that time. In particular $U=B$ Haar-almost surely.
\end{maintheorem}

\subsection*{Why the coefficient is an expected weight}
The positive-level direction walk is simple symmetric walk killed at
zero. Its physical duration also includes arithmetic holding times.
The main arithmetic estimate proves that a long excursion uses twice
its direction clock, with error negligible on the rare survival scale.
Each departure consequently has the same physical tail coefficient
$c=2/\sqrt\pi$. Chronological assembly gives $c\E M$ without assuming
independence of the lengths, offsets, or count.

From a fixed level $s$, the first return has weight $s$; later departures
have weight one. Their expected total weight is
$h_{\rm cap}(s,D)=s+\E_{s,D}M$. The level gains precisely what the
future count loses at a departure, making this capacity harmonic.
Section~\ref{harm:section} proves the state tail, conditioned-history
law, and triadic extension. Sections~\ref{sec:ordinary} and
\ref{sec:gao} transfer the coefficient to the original clocks and
counted populations. Section~\ref{sec:windows} gives two short
applications of the finite-horizon bound to published descent theorems.

\subsection*{The questions answered}
Gao states on p.~265 of \cite{Gao1993} that his proportions should tend
to one. Said Duran's September 2026 preprint \cite{Duran2026},
Theorems A and B, gave the first public proof of the qualitative
density-one conclusions. The author's companion papers
\cite{GaoDensityOne,GaoTao} give an independent proof with deficiency
of order $k^{-1/2}$. The theorem above determines the leading
deficiencies.

For integer $t\ge0$, set
\[
a(t)=\#\{0\le r<2^t:T^t(r)=T^t(r+1),\ o_t(r)=o_t(r+1)\}.
\]
The exact cylinder dictionary gives
\begin{equation}\label{eq:duran-law}
\PP(A>t)=1-a(t)/2^t=\PP(\sigma>t)
\end{equation}
for the balanced absorption time in Duran's Corollary C. These counts
are OEIS A398792 for $t\ge1$ \cite{OEIS398792}, and
\begin{equation}\label{eq:sequence-asymptotic}
a(t)=2^t(1-\kappa/\sqrt t+o(t^{-1/2})).
\end{equation}
This answers Remark 3.7, ``We do not know the true order of
$P(\sigma>t)$,'' and the separate quantitative request in Section 7(1).
Gao's p.~266 question about equal \emph{finite} heights and halving
counts is a separate open problem.

Terras--Lagarias parity cylinders supply the foundational correspondence
\cite{Terras1976,Lagarias1985}; symmetric-walk estimates are classical.
The stationary, killed-occupation, and arithmetic-clock arguments are
proved here. The geometry companion \cite{GeometryCompanion} uses these
laws to determine the shape of late pair and family synchronization.
None of the main rate proofs depends on that companion; an overview of
the series is \cite{Umbrella}. We give exact
probabilistic constants, but no numerical value of $\E M$ or effective
asymptotic onset. The integer statements retain possible individual
exceptions and assume no convergence to one.

% ---- exact/pair ----
\section{The canonical shortcut pair and its return structure}\label{sec:pair}

We use the maps and odd-count convention of Section~\ref{sec:intro}.
Unless another initial state is specified, the pair starts from
$(X_0,X_0+1)$ with $X_0$ Haar-distributed on $\Ztwo$.
All statements about fresh randomness
below are conditional on a finite revealed history.  We never condition on
survival for all time.

\subsection{Orientation, absorption, and binary cylinders}

At shortcut time $t$ retain an orientation and write the two current entries as
\begin{equation}\label{pair:affine}
 Y_t=3^{s_t}X_t+D_t,\qquad (s_t,D_t)\in\N\times\Z.
\end{equation}
Here $X_t$ is the entry chosen as base by the orientation; the original entries
continue to be tracked through every swap.  If $\epsilon=X_t\bmod2$, applying
$T$ to both entries gives the following canonical kernel.

\begin{proposition}[The canonical six-row kernel]\label{pair:kernel}
The state and base updates are
\[
\begin{array}{c|c|c}
\text{condition}&(s',D')&X'\\\hline
D\text{ even},\ \epsilon=0&(s,D/2)&T(X)\\
D\text{ even},\ \epsilon=1&(s,(3D+1-3^s)/2)&T(X)\\
D\text{ odd},\ s\ge1,\ \epsilon=0&(s+1,(3D+1)/2)&T(X)\\
D\text{ odd},\ s\ge1,\ \epsilon=1&(s-1,(D-3^{s-1})/2)&T(X)\\
D\text{ odd},\ s=0,\ \epsilon=0&(1,(3D+1)/2)&T(X)\\
D\text{ odd},\ s=0,\ \epsilon=1&(1,(1-3D)/2)&T(Y).
\end{array}
\]
The final row reverses the orientation.  All displayed offsets are integers,
and $(0,0)$ is absorbing.  Define the symbolic absorption time
\[
 A=\inf\{t\ge0:(s_t,D_t)=(0,0)\}.
\]
\end{proposition}
\begin{proof}
Since $3^s$ is odd, the entries have the same parity precisely when $D$ is
even.  Substituting the appropriate two branches of $T$ in
\eqref{pair:affine} gives the first four rows.  At level zero with odd $D$, an
even base gives the fifth row.  An odd base and even second entry instead give
$T(X)=3T(Y)+(1-3D)/2$; taking $T(Y)$ as the new base gives the final row.
\end{proof}

The boundary swap in this kernel is part of its definition.  In particular,
one must not replace it by the boundary update of a raw compact affine kernel
whose exponent or orientation is left unnormalized.  Compact-coordinate
representations used later must retain the appropriate deterministic adapter
to these six rows.

\begin{lemma}[Fresh fair labels and exact cylinders]\label{pair:cylinders}
Given any finite labelled history, the current base is Haar-distributed.
Thus its next parity is a fresh fair bit.  A word of $h$ labels determines
exactly one residue class of the initial base modulo $2^h$, and every such
class occurs.  In particular, with
\[
 S_h(d)=\PP_{(0,d)}(A>h),
\]
there are exactly $2^h S_h(d)$ nonabsorbed residues in a complete period
modulo $2^h$.
\end{lemma}
\begin{proof}
Write $X=2U+\epsilon$.  Given its parity, $U$ is Haar.  The next base is
$U$ when the base is even, $3U+2$ when it is odd without a swap, and
$U+(D+1)/2$ in the swapped row.  These affine maps have odd coefficient,
so each preserves Haar measure on $\Ztwo$.  This proves the conditional
statement by induction.  On residues the same induction says that a next
label chooses one of the two refinements of the current cylinder: the
new base coordinate still has odd coefficient in the remaining quotient.
Consequently each word has one cylinder of measure $2^{-h}$, and the
cylinders partition $\Ztwo$.  The assertion applies after absorption as
well, retaining the state $(0,0)$ and advancing the common base.
\end{proof}

For the two original entries $a,b$, the oriented level has a signed version:
it is $s_t$ when the current orientation agrees with the initial one and
$-s_t$ after reversal.  With a consistent choice of sign it equals
$o_t(b)-o_t(a)$.

\begin{lemma}[Exact balanced shortcut dictionary]\label{pair:dictionary}
For every integer initial pair $(a,a+d)$ and every $t\ge0$, canonical
absorption by time $t$ is equivalent to
\begin{equation}\label{pair:balancedshortcut}
 T^t(a)=T^t(a+d),\qquad o_t(a)=o_t(a+d).
\end{equation}
The same pointwise identity holds for pairs in $\Ztwo$; Haar measure
is used only for the cylinder counts.
For positive starts, absorption at a shortcut time $h$ produces a balanced
ordinary merger at time
\begin{equation}\label{pair:ordinaryclock}
 m=h+o_h(a)=h+o_h(a+d)\le2h,
\end{equation}
with exactly $h$ halving branches on each ordinary trajectory.
\end{lemma}
\begin{proof}
At every step the signed level changes by the difference of the two original
odd-branch indicators.  The six rows verify this identity by induction;
the last row reverses orientation exactly when this difference crosses
zero.  Starting from level zero, equality of the odd counts is therefore
equivalent to $s_t=0$.  At level zero, \eqref{pair:affine} becomes
$Y_t=X_t+D_t$, so equality of the endpoints is equivalent to $D_t=0$.
This proves absorption at time $t$.  Once absorbed, both entries and both
odd counts evolve identically, giving the equivalence with absorption by
$t$.  Each even shortcut step is one ordinary halving; each odd shortcut
step is an ordinary odd branch followed by a halving.  A completed grid of
$h$ shortcuts therefore has ordinary time $h+o_h$ and exactly $h$ halvings,
which proves \eqref{pair:ordinaryclock}.
\end{proof}

This is the endpoint-and-count observable in Duran's residue count
\cite[Corollary~C and Sections~2.2--2.4]{Duran2026}.  In particular, if $a(t)$
denotes that count for adjacent starts, the cylinder dictionary gives
\[
 \frac{a(t)}{2^t}=\PP(A\le t),
\]
and Duran's absorption variable $\sigma$ has the same distribution as $A$.
The preprint's use of an absolute offset can relabel boundary edges, so equality of
the observables does not assert equality of signed labelled state paths.
Likewise, \eqref{pair:ordinaryclock} is an implication about a completed
shortcut grid; the exact relation with the first ordinary clocks is proved in
Section~\ref{sec:ordinary}.

\subsection{The embedded walk and fractional moments}

When $s>0$, even offsets hold the level fixed and odd offsets move it up or
down by one with equal conditional probability.  Write
\begin{equation}\label{pair:valuation}
 e_j=v_2(3^j-1)=
 \begin{cases}1,&j\text{ odd},\\2+v_2(j),&j\text{ even}.
 \end{cases}
\end{equation}

\begin{lemma}[Holding and return]\label{pair:holding}
At the successive positive-level odd-offset transitions, the level is a
simple symmetric nearest-neighbor walk, stopped at zero.  At level $j$,
the conditional expected number of holding steps together with their
terminal odd-offset transition is at most $e_j+2$, uniformly in the
incoming offset.  Every holding segment is almost surely finite, and every
positive initial level returns to zero almost surely.
\end{lemma}
\begin{proof}
The fair directions follow from Lemma~\ref{pair:cylinders}.  For an even
offset with valuation below $e_j$, either branch decreases that valuation
by one.  A valuation above $e_j$ is sent to $e_j-1$ by the odd label,
whereas a valuation equal to $e_j$ is sent there by the even label.  A
failure keeps the valuation at least $e_j$.  This includes offset zero,
whose valuation is interpreted as infinity.  The waiting time for a
successful label is bounded by a geometric variable of mean two; there
are then at most $e_j-1$ holding steps and one terminal transition.
The formula for $e_j$ follows by factoring $3^j-1$ for odd $j$ and then
using repeated squaring for even $j$.

For the embedded walk killed at $\{0,H\}$, the probability of hitting
$H$ before zero from $s<H$ is $s/H$.  A nearest-neighbor walk avoiding
zero cannot stay forever in a fixed finite interval.  Sending $H$ to
infinity proves almost-sure return to zero.  Almost-sure finiteness of
the countably many holding segments transfers this conclusion to physical
shortcut time.
\end{proof}

\begin{lemma}[Uniform physical return domination]\label{pair:return-bound}\label{family:return-bound}
For the first return $R$ to level zero from any positive level $s$,
uniformly in the integer offset and conditional on the revealed entry past,
\begin{equation}\label{pair:return-envelope}
 \PP_{(s,D)}(R>t)\le\min(1,C_0s/\sqrt t),\qquad C_0=30.
\end{equation}
\end{lemma}
\begin{proof}
Kill the direction walk at $0,H$, $H>s$. Its expected visits to $j$ are
\[
 G_{0,H}(s,j)=2\min(s,j)(H-\max(s,j))/H\le2s.
\]
This solves the finite difference equation with zero boundary values.
The holding bound of Lemma~\ref{pair:holding} and
$\sum_{j<H}v_2(j)\le H-1$ give
$\sum_{j<H}(e_j+2)\le(9/2)(H-1)$. Thus the physical exit time has
expectation at most $9sH$, without a terminal holding. Gambler's ruin
and Markov give $\PP(R>t)\le s/H+9sH/t$.
For $t\ge(s+1)^2$, take $H=\lceil\sqrt t\rceil$; this is at most
$19s/\sqrt t$. For smaller $t$, $1\le2s/\sqrt t$. Both imply the
stated envelope. Each holding is averaged at its actual entry before
future bits, so the proof also holds conditionally on the finite past.
\end{proof}

Fix $0<\alpha<1$, and reserve $\kappa$ for the absorption constant rather
than the drift coefficient.  Set
\begin{equation}\label{pair:driftconstants}
 q=2^{-\alpha},\quad p=(3/2)^\alpha,\quad
 \chi_\alpha=(q+p)/2<1,\quad
 r=\frac{1-\sqrt{1-qp}}q\in(0,1),\quad a=rq<1.
\end{equation}
Strict convexity of $2^{-u}+(3/2)^u$ between its equal endpoint values
at $u=0,1$ gives $\chi_\alpha<1$.  The smaller root $r$ of
$qz^2-2z+p$ lies in $(0,1)$, and
$(qr+p/r)/2=1$.  Write
\[
 z=D/3^s,\qquad V(D)=|D|^\alpha,\qquad
 W(s,z)=r^s|z|^\alpha.
\]
In the order of the six kernel rows, the normalized offset updates are
\[
 z/2,\quad 3z/2+(3^{-s}-1)/2,\quad
 z/2+1/(2\cdot3^{s+1}),\quad 3z/2-1/2,\quad
 z/2+1/6,\quad -z/2+1/6.
\]
Using $|u+v|^\alpha\le |u|^\alpha+|v|^\alpha$, at positive even
offsets the conditional coefficient of $W$ is $\chi_\alpha$, and at
positive odd offsets it is $(qr+p/r)/2=1$.  The additive source is in
both cases bounded by $c_0r^{s-1}$, for a constant depending only on
$\alpha$.  At an odd level-zero departure it follows that
\begin{equation}\label{pair:departuremoment}
 \E[W(s_1,z_1)\mid D]\le aV(D)+r6^{-\alpha}.
\end{equation}

\begin{lemma}[Raw and coasted return moments]\label{pair:returnmoments}
If the process starts at $(s,D)$, $s>0$, and $\Delta$ is the raw offset at its
first return to zero, then
\[
 \E V(\Delta)\le r^s|D/3^s|^\alpha+b_\infty.
\]
For an excursion departing from an odd level-zero offset $d$, there is a
constant $b$ such that
\begin{equation}\label{pair:rawrecurrence}
 \E_d V(\Delta)\le aV(d)+b.
\end{equation}
This last bound also holds after the following level-zero even coast to an
odd offset or absorption.
\end{lemma}
\begin{proof}
The Green function for the embedded walk on the positive half-line killed
at zero is $G_s(j)=2\min(s,j)$.  Indeed the chance to hit $j$ before zero
is $s/j$ for $s\le j$ and one for $s\ge j$; starting from $j$, the chance
of returning to $j$ before zero after one step is $1-1/(2j)$.  Multiplying
the hitting probability by the resulting expected $2j$ visits gives the
formula.  If $\rho$ is the physical return time, holding bounds and Tonelli
give
\[
 \E\sum_{t<\rho}r^{s_t-1}
 \le2\sum_{j\ge1}\min(s,j)(e_j+2)r^{j-1}
 \le2\sum_{j\ge1}j(e_j+2)r^{j-1}<\infty.
\]
Stop the one-step drift inequality at $\rho\wedge n$ and use Fatou's lemma
on the nonnegative terminal moment.  Almost-sure return from
Lemma~\ref{pair:holding} proves the first bound.  From level one the source
sum has the sharper bound $2\sum_j(e_j+2)r^{j-1}$; together with
\eqref{pair:departuremoment} this proves \eqref{pair:rawrecurrence}.
At level zero a nonzero even offset changes to $D/2$ or $3D/2$, drops its
valuation by exactly one, and multiplies its conditional $V$ moment by
$\chi_\alpha$.  The coast therefore cannot increase expected $V$.
Offset zero is already absorbed.  No uniform-integrability claim is used
at the unbounded return time.
\end{proof}

\subsection{Compact trials and integrable boundary rewards}

A departure is a step whose initial state is $(0,d)$ with $d$ odd.  Let
$M$ count \emph{all} departures before absorption, including those within
any absorbing-word trial.  Let $d_i$ be the offset of departure $i$, and
$\Delta_i$ its raw offset on the next return to level zero.  Its physical
excursion length $R_i$ includes the departure step and every subsequent
positive-level step through that return, but excludes the following
level-zero coast.  Let $C_{\mathrm{coast}}$ count all level-zero even-offset coast steps.

\begin{lemma}[Finite boundary rewards]\label{pair:rewards}
From a level-zero offset $u$,
\begin{equation}\label{pair:rewardbound}
 \E_u M+\E_u\sum_{i\le M}V(\Delta_i)+\E_u C_{\mathrm{coast}}
 \le B_*(1+V(u))
\end{equation}
for a constant $B_*$ depending only on the fixed $\alpha$.  Absorption is
almost sure.
\end{lemma}
\begin{proof}
First remove an initial even coast.  It has exactly $v_2(u)$ steps for
$u\ne0$ and zero for $u=0$, and its endpoint has expected $V$ at most
$V(u)$.  The elementary inequality
\begin{equation}\label{pair:coastbound}
 v_2(u)\le\frac{V(u)}{\alpha\log2}\quad(u\ne0)
\end{equation}
bounds its expected length.

On odd level-zero offsets, follow each excursion and its coast to form
the boundary chain $Z_i$, with zero a cemetery state.  Its recurrence is
$\E[V(Z_{i+1})\mid\mathcal F_i]\le aV(Z_i)+b$.
Choose $R_0>b/(1-a)$, put $c=a+b/R_0<1$, and stop upon first entry
$\sigma$ into the finite set
\[
 \mathcal K=\{d\in\Z:d\text{ odd},\ V(d)\le R_0\}\cup\{0\}.
\]
For $q_i=\E[V(Z_i)\one_{\{\sigma>i\}}]$, killing on entry gives
\begin{equation}\label{pair:killedrecurrence}
 q_{i+1}\le aq_i+b\PP(\sigma>i)\le cq_i,
 \qquad \sum_{i\ge0}q_i\le\frac{\E V(Z_0)}{1-c},
\end{equation}
since $\PP(\sigma>i)\le q_i/R_0$.  Thus the expected number of departures
before entry is bounded by a constant times $\E V(Z_0)$.

Every integer boundary offset has a finite absorbing word.  Remove powers
of two with zero labels.  At an odd offset with $|d|>1$, the word $01$
for $d\equiv3\pmod4$, and the word $11$ for $d\equiv1\pmod4$, return
to level zero with offsets
\[
 (3d-1)/4\quad\hbox{and}\quad -(3d+1)/4,
\]
respectively.  In their indicated residue classes both have smaller
absolute value.  Repeat, removing any new powers of two, until reaching
$1$ or $-1$, whose absorbing words are $001$ and $101$.  Select one
absorbing word for each nonzero state in $\mathcal K$, and let $\ell$
be the largest selected length.  A trial reveals the selected number of
labels and follows its word when those labels agree; it succeeds with
conditional probability at least $p_0=2^{-\ell}$.

All prefixes of these finitely many trial words, and all possible states
after a failed trial of length at most $\ell$, form a finite set $G$.
Recover from a failed endpoint as follows: if its level is positive,
first return to zero, then take its even coast; run the boundary chain
until reaching $\mathcal K$ again.  The first-return moment bound and
\eqref{pair:killedrecurrence} give a uniform finite expected departure
count for this recovery, because the starting state lies in $G$.
The $\ell$ trial
steps themselves make at most $\ell$ departures.  This also counts a
departure whose excursion is unfinished at the failed endpoint: that
excursion is completed during the initial positive-level recovery,
which creates no further departure before its first return.

If $N_c$ counts the trials performed, conditional failure probabilities
give $\PP(N_c>j)\le(1-p_0)^j$ and $\E N_c\le p_0^{-1}$.
The event that round $j$ is performed is known at its start.  Conditional
expectation and Tonelli therefore bound the total expected departure
count by a uniform per-round bound times $\E N_c$.  The trial phase
starts in the finite set $\mathcal K$, so its bound is uniform in that
entry state.  Adding the pre-compact phase proves
\begin{equation}\label{pair:meanbound}
 m(u):=\E_u M\le B_0(1+V(u)).
\end{equation}
Each recovery returns almost surely, and the geometric trial bound proves
almost-sure absorption.

It remains to bound raw-return moments for the full count, without
distinguishing trial-word departures.  Put
\[
 Q_i=\E_u[V(\Delta_i)\one_{\{M\ge i\}}],\qquad Q_0=V(u).
\]
The initial coast and every subsequent coast do not increase expected
$V$.  Thus the expected $V$ of departure $i$, weighted by its occurrence,
is at most $Q_{i-1}$.  The raw-return recurrence gives, for every $i\ge1$,
\[
 Q_i\le aQ_{i-1}+b\PP_u(M\ge i).
\]
Summing the first $n$ inequalities and using $Q_i\ge0$ yields
\[
 (1-a)\sum_{i=1}^nQ_i\le aV(u)+b\E_u M.
\]
Monotone convergence and \eqref{pair:meanbound} therefore bound the
full sum of raw-return moments by a constant times $1+V(u)$.  Finally
\eqref{pair:coastbound} bounds every coast, including the initial one,
by this same moment reward.  This proves \eqref{pair:rewardbound}.
\end{proof}

\begin{lemma}[Boundary moments]\label{pair:moments}\label{pair:coasts}
For the canonical process from $(0,1)$,
\[
 1\le\E M<\infty,\qquad \E M^2<\infty,\qquad \E C_{\mathrm{coast}}<\infty.
\]
\end{lemma}
\begin{proof}
Time zero is a departure, so $M\ge1$.  The first moment follows from
Lemma~\ref{pair:rewards}, which also gives $\E C_{\mathrm{coast}}<\infty$. At the raw return of excursion $i$, the strong
Markov property and \eqref{pair:rewardbound} bound the conditional expected
number of future departures by $B_*(1+V(\Delta_i))$.  The pathwise identity
\[
 M^2=\sum_{i\le M}\bigl(1+2(M-i)\bigr)
\]
and Tonelli consequently give
\[
 \E M^2\le(1+2B_*)\E M+
 2B_*\E\sum_{i\le M}V(\Delta_i)<\infty.
\]
The reward proof counts excursions across failed trial endpoints as well
as departures inside trial words, so this is the second moment of the
full count $M$ defined above.
\end{proof}

\begin{proposition}[The physical excursion decomposition]\label{pair:decomposition}
Let $K_i$ be the number of positive-level odd-offset transitions in
excursion $i$. Almost surely,
\begin{equation}\label{pair:timesplit}
 A=C_{\mathrm{coast}}+\sum_{i\le M}R_i.
\end{equation}
Given the history just after departure $i$, $K_i$ has the law of the first
hit $K$ of zero by a simple symmetric walk started at one.  The event
that excursion $i$ occurs is known before its future walk is revealed.
\end{proposition}
\begin{proof}
Coasts at level zero cannot absorb a nonzero offset: their updates are
$D/2$ and $3D/2$.  A departure goes to level one.  Positive-level holdings
cannot absorb, so absorption occurs only on the last down-step of an
excursion, from level one with offset one.  The coasts and the complete
excursions therefore partition all steps through $A$, even when a compact
trial cuts across that partition.  The conditional walk law follows
from fresh fair labels and the almost-sure finite holding segments.
\end{proof}

\subsection{The uniform finite-horizon bound}

The following finite-horizon bound supplies the growing-block
applications. Its proof uses the boundary rewards and a deterministic
contraction period, before any exact asymptotic coefficient.

\begin{lemma}[Remaining absorption from a scalar state]\label{pair:remaining-tail}
For $s\ge0$, $D\in\Z$, put $z=D/3^s$. There is a constant $C$,
independent of this state, such that, for $t\ge1$,
\[
 \PP_{s,D}(A>t)\le
 \min\{1,C(s+1+|z|^\alpha)/\sqrt t\}.
\]
\end{lemma}
\begin{proof}
At $s>0$ first finish the positive-level interval $R_0$. Its tail is
bounded by $Cs/\sqrt t$ by Lemma~\ref{pair:return-bound}.
Its raw return offset $U$ satisfies
$\E|U|^\alpha\le r^s|z|^\alpha+b_\infty$ by
Lemma~\ref{pair:returnmoments}. At $s=0$ put $R_0=0$ and $U=D$.
Conditioning at this return, Lemma~\ref{pair:rewards} bounds both the
future departure mean and the coast mean by $C(1+|z|^\alpha)$.
A full departure interval consists of its first step and a positive-level
return from level one. Thus Lemma~\ref{pair:return-bound} gives,
uniformly at its actual entry,
\[
 \PP(R_i>x\mid\mathcal F_i)\le\min(1,C/\sqrt x),\qquad
 \E(R_i\wedge x\mid\mathcal F_i)\le2C\sqrt x.
\]
Here $\mathcal F_i$ is the predeparture information, on the event that
this departure occurs. Separate any interval longer than $x$, and
apply Markov to the sum of the remaining truncated lengths.
The occurrence event is known at entry, so conditional expectation
and Tonelli give
$\PP(\sum_{i\le M}R_i>x)\le C\E M/\sqrt x$.
Markov bounds the coast tail by its mean divided by $x$.
Use the union bound at threshold $t/3$ for these three components,
and enlarge the constant for bounded $t$. This proves the claim.
\end{proof}

\begin{theorem}[Uniform pair absorption]\label{pair:uniform-tail}\label{thm:pair}
There is an absolute constant $B\ge1$ such that, for every integer
level-zero offset $d$ and every real $t\ge1$,
\[
 \PP_{0,d}(A>t)\le
 \min\left\{1,B\sqrt{\frac{1+\log(1+|d|)}t}\right\}.
\]
Every fixed integer scalar state is absorbed almost surely.
\end{theorem}
\begin{proof}
Fix once and for all an $\alpha\in(0,1)$. Retain absorbed states at
$(0,0)$. For $z_h=D_h/3^{s_h}$, the six normalized rows and
subadditivity give
\[
 \E(|z_{h+1}|^\alpha\mid\mathcal F_h)
       \le\chi_\alpha|z_h|^\alpha+B_\alpha,
 \qquad
 \E|z_h|^\alpha\le\chi_\alpha^h|d|^\alpha+
                          \frac{B_\alpha}{1-\chi_\alpha}.
\]
At a positive odd-offset state the conditional squared-level increment
is one; at an even offset it is zero; at a level-zero departure it is
one. Hence $\E s_h^2\le h$ and $\E s_h\le\sqrt h$.
Choose the deterministic burn length
\[
 n_d=\left\lceil\frac{\log(1+|d|^\alpha)}{|\log\chi_\alpha|}
                         \right\rceil+1
       \le C(1+\log(1+|d|)).
\]
The normalized-offset moment at $n_d$ is bounded by an absolute
constant. Condition on the finitely many histories at this time and
apply Lemma~\ref{pair:remaining-tail}. For $t\ge2n_d$ this gives
\[
 \PP_{0,d}(A>t)
 \le \frac{C(\sqrt{n_d}+1)}{\sqrt{t-n_d}}
 \le C\sqrt{\frac{1+\log(1+|d|)}t}.
\]
For $t<2n_d$ enlarge the constant and use the probability bound one.
At $d=0$ absorption is immediate. Almost-sure absorption from arbitrary
integer states follows by their almost-sure positive-level return,
if needed, and Lemma~\ref{pair:rewards} at its integer boundary offset.
\end{proof}

This is an upper bound at every horizon and with a uniform absolute
constant. It supplies no uniform relative asymptotic. Its logarithmic
scale is necessary for such a bound: from $d=2^k$ the first $k$ steps
are nonzero boundary coasts, so $\PP(A>k)=1$.

% ---- exact/stationary ----
\section{Stationary rigidity of the raw kernel}
\label{st:section}

We first prove a rigidity statement for the compact raw kernel.  Its
hypothesis concerns only the level marginal; absolute continuity of the
joint law is neither assumed nor obtained by a preliminary smoothing
argument.  This distinction matters for the expected occupations used
below, whose finite approximants may be atomic.

Write $E=\Ztwo^2$, with raw coordinates $(s,D)$, and let $\lambda$ be
product normalized Haar measure.  The compact raw branches are
\[
\begin{array}{c|c|c}
 &F_0(s,D)&F_1(s,D)\\\hline
 D\text{ even}&(s,D/2)&(s,(3D+1-3^s)/2)\\
 D\text{ odd}&(s+1,(3D+1)/2)&(s-1,(D-3^{s-1})/2).
\end{array}
\]
These extend the positive ordinary-level rows to every $s\in\Ztwo$.
The ordinary canonical process swaps orientation at level zero; this
compact raw kernel instead retains the displayed signed decrement there.
The discrepancy is paid at the killed ordinary boundary in
\eqref{occ:boundary-telescoping}. We use the Markov operator
\[
 Pf=\tfrac12(f\circ F_0+f\circ F_1).
\]
All bits in this section are independent fair bits of this kernel.  The
continuous power is $3^{2j}=9^j$, $3^{2j+1}=3\cdot9^j$, where
$9^j=(1+8)^j$ is its binomial extension to $\Ztwo$.  It is a continuous
homomorphism from the additive group $\Ztwo$ into $\Ztwo^\times$.  In
particular,
\begin{equation}
 v_2(3^u-1)=
 \begin{cases}1,&u\text{ odd},\\2+v_2(u),&u\ne0\text{ even},\end{cases}
 \label{st:lte}
\end{equation}
and $s\mapsto3^s$ is injective.  For the even case, the identity
$v_2(9^j-1)=3+v_2(j)$ follows by successive squaring and continuity;
the odd case follows modulo four.  These observations also prove
injectivity by applying \eqref{st:lte} to the difference of two exponents.

\begin{lemma}[Haar odd-offset flux]\label{st:level-flux}
For every finite raw stationary measure $\mu$, the level measure
$\vartheta(B)=\mu\{s\in B,D\text{ odd}\}$ is $\vartheta(\Ztwo)$ times Haar.
\end{lemma}
\begin{proof}
If $\sigma$ is the level marginal, projection of the raw rows gives
\[
 \sigma=\sigma-\vartheta+
       \tfrac12(T_{+1}\vartheta+T_{-1}\vartheta),
\]
where $T_a$ translates by $a$.  On a nonconstant level character
$c(s)=\exp(2\pi i js/2^L)$ the translation average has multiplier
$\cos(2\pi j/2^L)<1$.  Hence $\vartheta(c)=0$.
Fourier inversion on each finite residue group proves the assertion.
\end{proof}

\begin{theorem}[Raw stationary uniqueness]\label{st:uniqueness}
Let $\mu$ be a probability measure on $E$ such that $\mu P=\mu$.  If its
level marginal $\mu_s$ is absolutely continuous with respect to normalized
Haar measure on $\Ztwo$, then $\mu=\lambda$.
\end{theorem}

The proof combines a coefficient escape argument, a finite word coding
argument, and an exact energy identity.  We give the components before
assembling them.  A qualitative mixing statement would not suffice for
this assembly: the energy decay and the partition precision must have
matching exponential scales.

\subsection{Escape of affine curve coefficients}
\label{st:affine-subsection}

For $A,B\in\Qtwo$, define the partial affine graph
\[
 G_{A,B}=\{(s,D)\in E:D=A+B3^s\}.
\]
It is Borel (indeed closed).  The use of coefficients in $\Qtwo$, rather
than only in $\Ztwo$, is necessary because a raw step divides by two.

\begin{lemma}[Affine graphs have zero stationary mass]\label{st:affine}
If $\mu P=\mu$ and $\mu_s$ is atomless, then every countable union of
sets $G_{A,B}$ has $\mu$-measure zero.
\end{lemma}
\begin{proof}
Put $X=3^s$.  Substitution of $D=A+BX$ into the raw branches gives the
following coefficient transformations, where $D'=A'+B'X'$:
\[
\begin{array}{c|c|c|c}
 D\bmod2&b&A'&B'\\\hline
 0&0&A/2&B/2\\
 0&1&(3A+1)/2&(3B-1)/2\\
 1&0&(3A+1)/2&B/2\\
 1&1&A/2&(3B-1)/2
\end{array}
\]
Thus, independently of the parity selector, the slope follows
\begin{equation}
 L_0(B)=B/2,\qquad L_1(B)=(3B-1)/2,
 \label{st:slope-kernel}
\end{equation}
each with probability $1/2$.

This coefficient kernel has no stationary probability on $\Qtwo$.  Once
$v_2(B)<0$, either map decreases the valuation by exactly one.  Before
exit from $\Ztwo$, at least one of the two maps exits at each step:
$L_1$ does so when $B$ is even, and $L_0$ when $B$ is odd.  Hence the
probability of remaining in $\Ztwo$ for $n$ further steps is at most
$2^{-n}$.  From every initial $B\in\Qtwo$ we therefore have
$|B_n|_2\longrightarrow\infty$ almost surely.  If a stationary
probability $\beta$ existed, dominated convergence, applied to
$\one_{\{|B_n|_2\le2^R\}}$, would give
$\beta\{|B|_2\le2^R\}=0$ for every integer $R$.  The countable union of
these balls is $\Qtwo$, a contradiction.

Close the given coefficient family under the four transformations in the
table.  The closure is countable, and its union $G$ is forward invariant.
Stationarity gives
\[
 \mu(G)=\mu(G)+\int_{G^c}P\one_G\,d\mu,
\]
so there is no incoming mass.  If $\mu(G)>0$, the normalized restriction
$\eta=\mu|_G/\mu(G)$ is stationary and has atomless level marginal.
We show that its slope labels contradict the preceding escape argument.

Distinct graphs $G_{A,B}$ and $G_{A',B'}$ intersect at
at most one level: if $B\ne B'$, their intersection fixes $3^s$, and
injectivity fixes $s$; if $B=B'$ and $A\ne A'$, they are disjoint.
All multiple-label points therefore form an $\eta$-null set.  Choosing
the first graph in an enumeration gives a Borel coefficient label,
unique almost everywhere.  Since $(F_b)_*\eta\le2\eta$, almost every
child also avoids the multiple-label set.  The table and forward closure
then identify its unique slope as $L_b(B)$.  Pushing $\eta$ forward by
the slope label produces the excluded stationary coefficient law.
Consequently $\mu(G)=0$, as required.
\end{proof}

\subsection{Word signatures and conditional nonatomicity}
\label{st:words-subsection}

Fix a word $w=(b_0,\ldots,b_{n-1})\in\{0,1\}^n$, and a parity history
$p=(p_0,\ldots,p_{n-1})$ realized on a source domain.  Set
\[
 q_w=\sum_{t<n}b_t,\qquad
 \ell_{w,p}=\sum_{t<n}p_t(1-2b_t).
\]
The endpoint level is $s+\ell_{w,p}$.  The multiplier of $D$ at step $t$
is $3/2$ precisely when $p_t\ne b_t$.  If $u$ counts these multipliers,
then $u-\ell_{w,p}=q_w$.  Define the coefficient signature recursively by
\begin{equation}
 C_0=0,\qquad
 C_{t+1}=\begin{cases}C_t,&b_t=0,\\3C_t+2^t,&b_t=1,\end{cases}
 \qquad C_w=C_n.
 \label{st:signature}
\end{equation}
On the chosen parity domain the endpoint offset has the form
\begin{equation}
 D_n=\frac{3^{q_w+\ell_{w,p}}D
                   -3^{\ell_{w,p}}C_w3^s+E_{w,p}}{2^n},
 \label{st:word-formula}
\end{equation}
where $E_{w,p}\in\Qtwo$ is independent of $(s,D)$.  One way to verify
all coefficients is to apply the word to the symbolic curve $D=A+B3^s$.
Its final slope, by the preceding table, is
$(3^{q_w}B-C_w)/2^n$, independently of $p$.  Multiplying by the endpoint
power $3^{s+\ell_{w,p}}$ gives the coefficient of $3^s$ in
\eqref{st:word-formula}; the coefficient of $D$ is $3^u/2^n$, and the
remaining constant is $E_{w,p}/2^n$.  Negative values of $\ell_{w,p}$
are harmless in $\Qtwo$.

\begin{lemma}[Injectivity of the coefficient signature]\label{st:signature-injective}
For a fixed length $n$, the pair $(q_w,C_w)$ determines $w$.
\end{lemma}
\begin{proof}
Write $w=(b_0,w')$.  Expanding \eqref{st:signature} gives
\[
 q_{w'}=q_w-b_0,\qquad
 C_w=b_0\,3^{q_{w'}}+2C_{w'}.
\]
Thus $b_0$ is the parity of the ordinary integer $C_w$, and the signature
of $w'$ is recovered by subtraction and division by two.  Induction on
the fixed length recovers every bit.
\end{proof}

\begin{lemma}[Free endpoint trees]\label{st:free-tree}
Under the hypotheses of Lemma~\ref{st:affine}, for $\mu$-almost every
source point and every $n\ge1$, the $2^n$ literal length-$n$ bit words
have distinct endpoints.
\end{lemma}
\begin{proof}
Take distinct equal-length words, together with two of their possible
parity histories.  If their $\ell$ values differ, their endpoint levels
differ by a nonzero ordinary integer in $\Ztwo$.  If the $\ell$ values
agree, subtract \eqref{st:word-formula}.  When their $q$ values differ,
the coefficient of $D$ is nonzero; endpoint equality defines one affine
graph $D=A+B3^s$.  Its measure is zero by Lemma~\ref{st:affine}.  When
the $q$ values agree, their $C$ values differ by
Lemma~\ref{st:signature-injective}.  The endpoint equality then fixes
$3^s$ to at most one value and so fixes at most one level, also a null
set.  There are finitely many parity histories at each length, and only
countably many lengths.  Removing the union of these null collision sets
proves the simultaneous statement.
\end{proof}

\begin{lemma}[Graph-null disintegration]\label{st:nonatomic}
If $\mu P=\mu$ and its level marginal $\sigma$ is atomless, then its
conditional offset law $\gamma_s$ is nonatomic for $\sigma$-almost every
$s$.
\end{lemma}
\begin{proof}
Let $G$ be any Borel set with at most $M$ offsets in its fibre at
each level.  A length-$n$ word can end only at one of the $2n+1$ levels
$s-n,\ldots,s+n$.  On the full-measure set of Lemma~\ref{st:free-tree},
at most $M(2n+1)$ distinct endpoints can belong to $G$.  Each literal
word has probability $2^{-n}$, and therefore
\[
 P^n\one_G(z)\le M(2n+1)2^{-n}\quad\text{for $\mu$-almost every }z.
\]
Stationarity and integration imply $\mu(G)=0$.  No regularity of the
finite fibre sections is required.

Disintegrate $\mu$ over $\sigma$ as conditional probabilities $\gamma_s$.
The singleton-mass function $(s,D)\mapsto\gamma_s\{D\}$ is measurable:
evaluate the kernel $(s,D)\mapsto\gamma_s$ on the section of the Borel
diagonal.  For $k\ge1$, the Borel set
\[
 E_k=\{(s,D):\gamma_s\{D\}\ge1/k\}
\]
has at most $k$ points in every level fibre, and hence $\mu(E_k)=0$.
Their union is the positive-singleton-mass set.  Disintegration shows
that, for $\sigma$-almost every $s$, this set has $\gamma_s$-measure
zero.  A positive atom of $\gamma_s$ would itself give that set positive
measure.  Thus $\gamma_s$ is nonatomic almost everywhere.
\end{proof}

\subsection{The critical collision estimate}
\label{st:collision-subsection}

Let $\mathcal C_q$ be the joint partition by $(s,D)$ modulo $2^q$.
The following estimate uses equal precision in both coordinates and
the conditional offset law given the full level.

\begin{lemma}[Critical little-o collision]\label{st:collision}
Let $\alpha$ be a finite measure on $E$ with level marginal $\sigma$
bounded above by $K$ times Haar measure.  Suppose its conditional offset
laws are nonatomic almost everywhere.  Then
\begin{equation}
 2^q\sum_{C\in\mathcal C_q}\alpha(C)^2\longrightarrow0.
 \label{st:critical}
\end{equation}
\end{lemma}
\begin{proof}
The zero measure is trivial.  Disintegrate a nonzero $\alpha$ as
$\sigma(ds)\gamma_s(dD)$ with probability conditionals.  Let $L$ be a
precision-$q$ level cell and $C_d=\{D\equiv d\pmod{2^q}\}$ an offset
cell.  Integral Cauchy--Schwarz gives
\[
 \alpha(L\times C_d)^2
 =\left(\int_L\gamma_s(C_d)\,d\sigma(s)\right)^2
 \le \sigma(L)\int_L\gamma_s(C_d)^2\,d\sigma(s).
\]
Since $\sigma(L)\le K2^{-q}$, summing over the two partitions yields
\begin{equation}
 2^q\sum_{C\in\mathcal C_q}\alpha(C)^2
 \le K\int\sum_{d\bmod2^q}\gamma_s(C_d)^2\,d\sigma(s).
 \label{st:collision-bound}
\end{equation}
For each $s$, the inner sum is the probability that two independent
samples from $\gamma_s$ agree modulo $2^q$.  These events decrease to
equality of the samples, whose probability is
$\int\gamma_s\{D\}\,d\gamma_s(D)$.  It is zero for almost every $s$
by conditional nonatomicity.  The inner sum is measurable and bounded
by one.  Bounded convergence under the finite measure $\sigma$ proves
\eqref{st:critical} directly as $q\to\infty$.
\end{proof}

For a stationary law with level density $w\in L^1$, the finite restriction
$\alpha_K=\mu|_{\{w(s)\le K\}}$ has bounded level density and inherits
nonatomic conditional offsets from Lemma~\ref{st:nonatomic}.  Hence it
satisfies \eqref{st:critical}.  Its omitted mass tends to zero with $K$.
We do not assert that this restriction is stationary.

\subsection{High level bits and exact Haar energy}
\label{st:energy-subsection}

Each branch $F_b$ preserves $\lambda$.  To check this branch by branch,
write $D=2u$ or $D=2u+1$ on the two parity halves.  Conditional on $s$,
$u$ is Haar; the output offset is an odd affine bijection of $u$, and the
output level is $s$ or a fixed translate of $s$.  Each half therefore
pushes its mass to $\lambda/2$.  The operators $U_bf=f\circ F_b$ are
$L^2(\lambda)$ isometries, and $P=(U_0+U_1)/2$ is an $L^2$ contraction.

\begin{proposition}[Exact energy halving]\label{st:energy}
Fix a bounded level cylinder $c$ of precision $L\ge0$, an odd integer
$k$, and $N\ge1$.  Set
\[
 f(s,D)=c(s)\exp(2\pi i kD/2^N),\qquad
 n_0=\max(3,L+3-N).
\]
For $n\ge n_0$,
\begin{equation}
 \|P^nf\|_{L^2(\lambda)}^2
       =\tfrac12\|P^{n-1}f\|_{L^2(\lambda)}^2.
 \label{st:energy-halving}
\end{equation}
In particular $\|P^nf\|_2^2\le C_f2^{-n}$ for all $n\ge0$, with a
constant depending on this fixed test.  If
$q_0=\max(L,N,1)$, then $P^nf$ has joint precision at most $q=n+q_0$.
\end{proposition}
\begin{proof}
For $n\ge n_0$, put
\[
 h=n+N-3\ge1,\qquad \Delta=2^h,\qquad r=n+N-1.
\]
By \eqref{st:lte}, $v_2(3^\Delta-1)=r$.  Fix a literal length-$n$
word.  Compare its paths from $(s,D)$ and $(s+\Delta,D)$, and let
$\delta_t$ be the translated path's offset minus the original path's
offset.  Throughout these paths the level difference is $\Delta$ and
\begin{equation}
 v_2(\delta_t)\ge r-t\quad(0\le t\le n),
 \label{st:highbit-valuation}
\end{equation}
where $v_2(0)=\infty$.  Normalize the difference by
$\eta_t=2^t\delta_t/(3^\Delta-1)$.  Whenever the input offsets have
the same parity, the raw rows give equal level increments and
\[
 \eta_{t+1}=a_t\eta_t-b_t2^t3^{s_t-e_t},
 \qquad a_t\in\{1,3\},\quad e_t\in\{0,1\}.
\]
Here $e_t=0$ for the even row and $e_t=1$ for the odd row; the exponent
is an arbitrary element of $\Ztwo$.  Starting from $\eta_0=0$, induction
keeps $\eta_t$ integral and proves \eqref{st:highbit-valuation}.  At each
input time $t<n$ its lower bound is at least $N\ge1$, so the offsets
have the same parity and the induction is valid through the last step.
The normalized recurrence also gives $\eta_n\equiv b_0\pmod2$:
its first injection has parity $b_0$, and all later injections are even.
Since $(3^\Delta-1)/2^n$ has valuation $N-1$, this means
$\delta_n\equiv b_0 2^{N-1}\pmod{2^N}$.  The endpoint level difference
$\Delta$ is divisible by $2^L$, and therefore
\begin{equation}
 f(F_w(s+\Delta,D))=(-1)^{b_0}f(F_w(s,D)).
 \label{st:highbit-sign}
\end{equation}
No branch probabilities or parity histories were conditioned or changed
in this comparison.

Let $A_n=U_0P^{n-1}f$ and $B_n=U_1P^{n-1}f$.  Averaging
\eqref{st:highbit-sign} over suffix words shows that level translation
by $\Delta$ fixes $A_n$ and negates $B_n$.  Since that translation
preserves Haar measure,
$\int A_n\overline{B_n}\,d\lambda=0$.  Expanding
$\|(A_n+B_n)/2\|_2^2$ and using the branch isometries gives
\eqref{st:energy-halving}.  Iteration and contraction on the finitely
many earlier iterates give the asserted bound, for example
$C_f=2^{n_0-1}\|f\|_2^2$.

Finally a function of joint precision $q\ge1$ pulls back by either raw
branch to a function of precision at most $q+1$.  The additional bit
fixes the parity and supplies the bit consumed by division by two; the
powers of three in the affine constants are determined at the required
precision by those same input level bits.  Hence induction starting at
$q_0$ proves the precision assertion.
\end{proof}

\begin{proof}[Proof of Theorem~\ref{st:uniqueness}]
The absolutely continuous level marginal is atomless, so the preceding
lemmas apply.  Fix a nonzero-offset Fourier term $f$ as in
Proposition~\ref{st:energy}.  A precision-$q$ joint cell has Haar mass
$2^{-2q}$.  For any finite measure $\alpha$, Cauchy--Schwarz on these
cells, with $q=n+q_0$, gives
\begin{align}
 |\alpha(P^nf)|
 &\le\|P^nf\|_{L^2(\lambda)}
       \left(2^{2q}\sum_{C\in\mathcal C_q}\alpha(C)^2\right)^{1/2}
 \notag\\
 &\le C'_f\left(2^q\sum_{C\in\mathcal C_q}\alpha(C)^2\right)^{1/2}.
 \label{st:energy-collision}
\end{align}
The constant is independent of $n$ because $q-n=q_0$ is fixed.
Apply this to $\alpha_K$ above.  Lemma~\ref{st:collision} makes it tend
to zero for each fixed $K$.  The omitted measure contributes at most
$\|f\|_\infty\mu\{w>K\}$, since $P$ is a sup-norm contraction.
Stationarity of the full $\mu$ gives $\mu(f)=\mu(P^nf)$.  First taking
$n\to\infty$, then $K\to\infty$, proves $\mu(f)=0$.

For a nonconstant level character $c$, the preceding offset result gives
$\mu(c(-1)^D)=0$.  The odd-offset flux $\vartheta$ of $\mu$ satisfies,
by Lemma~\ref{st:level-flux},
$0=\vartheta(c)=\mu(c)/2$, identifying the level marginal as Haar.
Every nonconstant joint
Fourier coefficient of $\mu$ vanishes.  Fourier inversion on each finite
joint residue group determines its cylinder masses as those of
$\lambda$; these cylinders determine the Borel probability law.  Thus
$\mu=\lambda$.
\end{proof}

\begin{theorem}[Full raw stationary classification]\label{st:classification}
Every probability measure $\mu$ on $E$ with $\mu P=\mu$ has the form
\[
 \mu=c\,\delta_{(0,0)}+(1-c)\lambda,\qquad 0\le c\le1.
\]
In particular, zero mass at the origin suffices for $\mu=\lambda$.
\end{theorem}
\begin{proof}
Fix a Haar-null Borel level set $B$.  Its odd-offset section has zero
$\mu$-mass by Lemma~\ref{st:level-flux}.  Branch domination gives
$(F_0^n)_*\mu\le2^n\mu$.  Every point with $s\in B$ and
$v_2(D)=n<\infty$ reaches that null odd-offset section after exactly
$n$ bit-zero steps: until then the level remains $s$ and the offset is
divided by two.  Thus $\mu\{s\in B,D\ne0\}=0$, by a countable union
over $n\ge0$.  The bit-one image of $(s,0)$ is
$(s,(1-3^s)/2)$.  For $s\ne0$ its offset is nonzero by
\eqref{st:lte}; domination by $2\mu$ therefore also gives
$\mu\{s\in B\setminus\{0\},D=0\}=0$.
Consequently the level marginal of $\mu$ away from the origin is
absolutely continuous.

Put $c=\mu\{(0,0)\}$.  Both branches fix the origin, so
$\mu-c\delta_{(0,0)}$ is a positive stationary measure.  Its level
marginal is absolutely continuous by the null-set argument just proved.
If $c<1$, normalize it and apply Theorem~\ref{st:uniqueness}.
Both the origin and product Haar are stationary, which completes the
classification.
\end{proof}

\begin{remark}\label{st:precision-scope}
The input test in Proposition~\ref{st:energy} has fixed precision.  Its
\emph{pullback} precision $n+q_0$ moves with $n$, and the matching
little-o collision estimate is proved at that moving precision.  This
does not make the conclusion uniform over a family of input tests whose
precision grows.  The stationary control $\delta_{(0,0)}$ also explains
why an absorbing component must be excluded: both raw branches fix
$(0,0)$, and its level law is atomic.  The theorem makes no pointwise assertion about free trees
from every ordinary integer start.
\end{remark}

\section{Killed expected occupation and fixed cylinders}
\label{occ:section}

Let $\mathcal A=\{s\text{ odd},D\text{ odd}\}\subset E$, and identify it with
$\mathcal X=\Ztwo\times\Ztwo^\times$ by $s=2J+1$.  Let $\pi$ be Haar measure
in $J$ times normalized Haar measure on odd $D$; thus
$\pi=4\lambda|_{\mathcal A}$ under this chart.  Let $\tau$ be the first subsequent
raw return to $\mathcal A$.  For $J\ne0,-1$ the raw return is almost surely finite,
as proved below; write $Q$ for its law there.  On the two exceptional
fibres $J=0,-1$, give $Q$ the identity extension.  This defines a Borel
kernel.  We will neither assume global Feller continuity of this
extension nor replace finite literal returns there by the identity.

The canonical killed epoch process has states $Z_v=(J_v,D_v)$ for
$v<\tau_{\mathrm e}$.  We allow any initial probability law $\xi$ on
ordinary $J\ge0$ and odd $D\in\Ztwo$, including dependent coordinates.
The original departure starts at $J_0=0$ with an odd-offset law $\beta$;
it is the special case $\xi=\delta_0\otimes\beta$.
The level transition probabilities are
\begin{equation}
 \begin{array}{c|ccc}
 j\ge1&j-1:1/4&j:1/2&j+1:1/4\\
 j=0&\dagger:1/2&0:1/4&1:1/4.
 \end{array}
 \label{occ:killed-level}
\end{equation}
The cemetery state $\dagger$ is assigned test value zero.  These
probabilities follow from the two fresh direction bits at their actual
entry times, with intervening horizontal holdings.  The special killed
row at zero retains the ordinary departure/arrival convention.
Away from it the raw return law and the killed epoch transition agree.

Define the probability measures
\begin{equation}
 a_m(\xi)=\E_\xi\min(\tau_{\mathrm e},m),\qquad
 \mu_{m,\xi}=\frac1{a_m(\xi)}\E_\xi
       \sum_{v<\min(\tau_{\mathrm e},m)}\delta_{Z_v}.
 \label{occ:occupation}
\end{equation}
The denominator depends only on the initial level law.  Every visit
has its actual survival weight: this is the expected total occupation
normalization.  For $\xi=\delta_0\otimes\beta$, write
$a_m=a_m(\xi)$ and $\mu_{m,\beta}=\mu_{m,\xi}$; then $a_m$ is
independent of $\beta$.

\subsection{Green bounds and periodic correctors}
\label{occ:level-subsection}

\begin{lemma}[Killed level occupation]\label{occ:green}
For \eqref{occ:killed-level}, starting at any ordinary $h\ge0$,
$\tau_{\mathrm e}<\infty$ almost surely.  The total Green kernel is
\begin{equation}
 G(h,j):=\E_h\#\{v<\tau_{\mathrm e}:J_v=j\}
       =2(2\min(h,j)+1)\le2(2j+1).
 \label{occ:green-kernel}
\end{equation}
In particular $G(h,0)=2$ for every $h$.  The survival probability is
\begin{equation}
 \PP_h(\tau_{\mathrm e}>v)=4^{-v}
       \sum_{k=-h}^{h}\binom{2v}{v+k},
 \label{occ:survival-start}
\end{equation}
where binomial coefficients outside their usual range are zero.
Writing $a_m$ for the zero-start denominator,
\begin{equation}
 a_m=2m\binom{2m}{m}4^{-m}\longrightarrow\infty,\qquad
 a_m(\xi)\ge a_m
 \label{occ:normalizer}
\end{equation}
for every initial law $\xi$.
\end{lemma}
\begin{proof}
Use an ordinary simple symmetric walk started at $2h$, killed on reaching
$-1$, and sample every two steps.  Its even-time surviving state $2J_v$
has exactly \eqref{occ:killed-level}.  Reflection at the first hit of
$-1$ gives the endpoint count at $2j$ as
\[
 \binom{2v}{v+j-h}-\binom{2v}{v+j+h+1}.
\]
Summing over $j\ge0$ telescopes to \eqref{occ:survival-start}.
Each binomial term is at most the central one, so for fixed $h$ the
survival probability is at most $(2h+1)c_v\to0$, where
$c_v=\binom{2v}{v}4^{-v}$.  This proves almost-sure killing.

For fixed target $j$, put $V_j(h)=2(2\min(h,j)+1)$ and
$V_j(\dagger)=0$.  For $j\ge1$ its interior slope drops from four to
zero at $j$; the zero row has $V_j(0)=2,V_j(1)=6$.
For $j=0$, $V_0\equiv2$ on surviving states.  These substitutions give
\[
 (I-P_{\mathrm{lev}})V_j(h)=\one_{\{h=j\}}.
\]
Here $P_{\mathrm{lev}}$ is the killed level kernel.  Stopped telescoping
at $\tau_{\mathrm e}\wedge m$ gives
\[
 \E_h\sum_{v<\tau_{\mathrm e}\wedge m}\one_{\{J_v=j\}}
 =V_j(h)-\E_h V_j(J_{\tau_{\mathrm e}\wedge m}).
\]
The bounded potential vanishes at the almost-sure killing time, so
bounded and monotone convergence prove \eqref{occ:green-kernel}.

The survival sum also includes the central term, giving
$\PP_h(\tau_{\mathrm e}>v)\ge c_v$.  The recursion
$c_{v+1}=c_v(2v+1)/(2v+2)$ gives
$\sum_{v<m}c_v=2mc_m$, by induction or telescoping.
Since $c_m\sim(\pi m)^{-1/2}$, this proves divergence and, by averaging
the survival lower bound over the initial law, $a_m(\xi)\ge a_m$.
\end{proof}

\begin{lemma}[Uniform Haar level limit]\label{occ:level-haar}
For every fixed continuous level test $r$, the $J$ marginal of
$\mu_{m,\xi}$ converges to Haar, uniformly over all initial laws $\xi$.
For every fixed ordinary $C\ge0$,
\begin{equation}
 \mu_{m,\xi}\{0\le J\le C\}\le\frac{2(C+1)^2}{a_m(\xi)}.
 \label{occ:finite-levels}
\end{equation}
For the departure specialization $\xi=\delta_0\otimes\beta$, the sharper
bound is $\mu_{m,\beta}\{J\le C\}\le2(C+1)/a_m$.
\end{lemma}
\begin{proof}
For a mean-zero real test $r$ on the cyclic group $\Z/q\Z$, solve the
finite periodic Poisson equation
\[
 [g(j+1)-2g(j)+g(j-1)]/4=r(j)
\]
with bounded $q$-periodic $g$.  The discrete Fourier Laplacian has kernel
exactly the constants, so the mean-zero assumption is precisely the
solvability condition.  At interior states the killed drift of $g$ is
$r$.  At zero it is $r(0)+c$, where
$c=-(g(0)+g(-1))/4$, with $g(\dagger)=0$.  Bounded stopping at
$\min(\tau_{\mathrm e},m)$ and the tower property therefore give
\[
 \left|\E_\xi\sum_{v<\min(\tau_{\mathrm e},m)}r(J_v)\right|
 \le2\|g\|_\infty+|c|\E_\xi\#\{v<\tau_{\mathrm e}:J_v=0\}
 \le3\|g\|_\infty.
\]
Division by $a_m(\xi)\ge a_m$ proves equidistribution for each fixed modulus,
uniformly in $\xi$.  In particular each precision-$a$ level
cylinder has limiting mass $2^{-a}$.  Level cylinders are uniformly dense
in continuous functions on $\Ztwo$, so the claim extends to those tests.
Summing the target-level bound $G(h,j)\le2(2j+1)$ over $0\le j\le C$
gives \eqref{occ:finite-levels}, since $\sum_{j=0}^C(2j+1)=(C+1)^2$.
At $h=0$ the Green mass at each level is two, giving the sharper
specialized bound.
\end{proof}

\subsection{Return times and exceptional fibres}
\label{occ:return-subsection}

\begin{lemma}[Offset-uniform holding and return bounds]\label{occ:holding}
At a nonzero raw level $s$, the expected horizontal holding time from an
even offset until an odd offset is at most $v_2(3^s-1)+1$, uniformly in
the initial offset.  At an odd offset it is zero.  For $r\ge2$, put
\[
 G_r=\{J\not\equiv0,-1\pmod{2^r}\}.
\]
Then $\E_z\tau\le r+7$ on $G_r$.  If an epoch law has Haar $J$
marginal, its mean raw return time is at most nine.
\end{lemma}
\begin{proof}
Write $e=v_2(3^s-1)$.  During even-offset holding the two choices are
$D/2$ and $(3D+1-3^s)/2$.  If $1\le v_2(D)<e$, either choice decreases
the valuation by one.  If $v_2(D)\ge e$, at least one choice lands at
valuation $e-1$: use bit zero when $v_2(D)=e$, and bit one when
$v_2(D)>e$, including $D=0$.  Until this happens the conditional chance
of that event is at least $1/2$, so its waiting time has mean at most two.
Thereafter at most $e-1$ further steps reach an odd offset.  The total
mean is at most $e+1$.  This proves the bound without assumptions about
the entry-offset distribution.

An epoch return outside $J=0,-1$ consists of two direction steps, one
even-level holding and one odd-level holding.  The intermediate even
level is $2J$ or $2(J+1)$ and is nonzero; the later odd level is also
nonzero.  On $G_r$ its even-level valuation is at most $r$.
Equation~\eqref{st:lte} gives even holding mean at most $r+3$ and odd
holding mean at most two, hence the budget $r+7$.  In particular the
return is almost surely finite at every nonexceptional epoch point.

For Haar $J$, the first fair direction enters even level $2U$ with
$U=J$ or $J+1$.  Each of those unconditioned alternatives is Haar,
regardless of the possible dependence of $D$ on $J$, since the direction
bit is fresh and fair before it is used.  Haar $U$ satisfies
$\E v_2(U)=\sum_{k\ge1}2^{-k}=1$, and $U=0$ has zero probability.
The even holding mean is at most $4+v_2(U)$, hence at most five on
average.  Add two direction steps and odd holding mean at most two to
obtain $\E\tau\le9$.
\end{proof}

\begin{lemma}[Actual occupation limits are stationary]\label{occ:stationary-limit}
Every weak limit of $\mu_{m_\ell,\xi_\ell}$, where $m_\ell\to\infty$
and the initial laws $\xi_\ell$ on ordinary nonnegative levels and odd
offsets may vary, has Haar $J$ marginal and is $Q$-stationary.
\end{lemma}
\begin{proof}
Compactness of $\mathcal X$ gives subsequential limits.  Lemma~\ref{occ:level-haar}
proves the Haar marginal along any such subsequence, and hence zero mass
on the exact fibres $J=0,-1$.

Let $f$ be a continuous epoch cylinder and $B=\|f\|_\infty$.  Define
\begin{equation}
 Q_Hf(z)=\E_z[f(X_\tau)\one_{\{\tau\le H\}}]
 \label{occ:truncated-return}
\end{equation}
by literal raw returns at \emph{all} epoch points, including the two
exceptional fibres.  Give value zero only to words with no return by
$H$; the raw endpoint $X_\tau\in\mathcal A$ is identified with its epoch chart
when evaluating $f$.  The raw maps and the parity domains are continuous, and the return
event through a fixed number of steps is clopen.  The finite sum over
words in \eqref{occ:truncated-return} is therefore a continuous finite
cylinder.  In particular it is not the truncation obtained by inserting
the identity extension on exceptional fibres.

On $G_r$, Lemma~\ref{occ:holding} and Markov's inequality yield
\begin{equation}
 |Qf-Q_Hf|\le B(r+7)/H;
 \qquad |Qf-Q_Hf|\le2B\text{ on }G_r^c.
 \label{occ:return-approximation}
\end{equation}
The second bound includes the possible discrepancy between literal
finite returns and the exceptional-fibre identity extension.  Uniform
level convergence gives
$\mu_{m,\xi}(G_r^c)\to2^{1-r}$ at fixed $r$, uniformly in $\xi$,
and any limit $\nu$ gives this set exactly mass $2^{1-r}$.

Let $P_{\mathrm{kill}}$ be the actual one-epoch killed transition, with
cemetery value zero.  Bounded stopped telescoping gives
\[
 |\mu_{m,\xi}(P_{\mathrm{kill}}f-f)|\le2B/a_m(\xi).
\]
At ordinary $J\ge1$, this transition equals $Q$.  It can differ only
at $J=0$, whose expected visits are at most two.  The extra cost is at
most $2B\cdot2/a_m(\xi)$.  Thus the complete boundary payment is
\begin{equation}
 |\mu_{m,\xi}(Qf-f)|\le6B/a_m(\xi)\longrightarrow0.
 \label{occ:boundary-telescoping}
\end{equation}
This pays the ordinary killed boundary before passage to a compact limit.

Along a weakly convergent subsequence, pass $Q_Hf-f$ to the limit at
fixed $r,H$.  The two approximation errors in
\eqref{occ:return-approximation}, one before and one after the limit,
are bounded in the limit by
\[
 2B(r+7)/H+4B\,2^{1-r}.
\]
Use \eqref{occ:boundary-telescoping}, then take $H\to\infty$ and next
$r\to\infty$.  This proves $\nu(Qf-f)=0$.  Finite cylinder tests
determine probability laws on $\mathcal X$, so $\nu Q=\nu$.  All estimates used
here are uniform over all initial laws $\xi$.
\end{proof}

\subsection{The raw occupation lift}
\label{occ:lift-subsection}

\begin{theorem}[Finite-mean literal-return uniqueness]\label{occ:epoch-uniqueness}
Let $\nu$ be a probability law on $\mathcal X$ whose literal first
subsequent raw return law is again $\nu$, and suppose
$M_\nu=\E_\nu\tau<\infty$.  Then $\nu=\pi$ and $M_\nu=4$.
In particular, these conclusions hold if $\nu Q=\nu$ and $\nu_J$ is Haar.
\end{theorem}
\begin{proof}
Finite mean implies an almost-sure return.  Define the raw probability
\begin{equation}
 \rho(B)=\frac1{M_\nu}\int_{\mathcal X}\E_z
        \sum_{0\le t<\tau}\one_B(X_t)\,d\nu(z).
 \label{occ:lift}
\end{equation}
For a bounded raw test $g$, stopped telescoping gives
\[
 M_\nu\rho(Pg-g)
   =\int_{\mathcal X}\E_z[g(X_\tau)-g(X_0)]\,d\nu(z)=0.
\]
First stop at $\tau\wedge H$; the summands are bounded by
$2\|g\|_\infty$ and their number has finite mean, so dominated
convergence justifies passage to $\tau$.  The literal endpoint law equals
the initial law $\nu$, proving $\rho P=\rho$.

A pre-return interval cannot visit the absorbing origin: after doing so
it would never return to $\mathcal A$.  Since return is almost sure,
$\rho\{(0,0)\}=0$.  Theorem~\ref{st:classification} therefore gives
$\rho=\lambda$, with no level-density construction.
A return interval contains exactly one visit to $\mathcal A$ before
its endpoint, its starting point.  Thus, as measures on $\mathcal A$,
\[
 \rho|_{\mathcal A}=\nu/M_\nu.
\]
Since $\lambda(\mathcal A)=1/4$, total masses give $M_\nu=4$ and
restriction gives $\nu=4\lambda|_{\mathcal A}=\pi$.

For the final assertion, Haar $J$ gives zero mass to $J=0,-1$, so $Q$
agrees $\nu$-almost everywhere with the literal return law.
Lemma~\ref{occ:holding} supplies almost-sure return and $M_\nu\le9$.
The preceding argument applies.
\end{proof}

\begin{theorem}[Uniform fixed-test expected occupation]\label{occ:uniform-convergence}
For every fixed continuous $f$ on $\mathcal X$,
\begin{equation}
 \sup_\xi|\mu_{m,\xi}(f)-\pi(f)|\longrightarrow0,
 \label{occ:uniform-limit}
\end{equation}
where the supremum ranges over every initial law $\xi$ on ordinary
nonnegative levels and odd offsets, with denominator $a_m(\xi)$.
In particular every fixed ordinary departure whose first epoch has $J_0=0$ and odd offset satisfies this conclusion with its actual
first-epoch offset law.
\end{theorem}
\begin{proof}
If \eqref{occ:uniform-limit} failed for a fixed test, there would be
$\epsilon>0$, $m_\ell\to\infty$, and $\xi_\ell$ for which the
absolute difference is at least $\epsilon$.  Compactness gives a weakly
convergent subsequence of these probability measures.  Its limit is
$Q$-stationary with Haar level by Lemma~\ref{occ:stationary-limit},
and equals $\pi$ by Theorem~\ref{occ:epoch-uniqueness}.  Weak convergence
for the fixed continuous test contradicts the difference.  The law at
the first epoch may itself be random after the departure's initial
holding; its law is the allowed special case $\xi=\delta_0\otimes\beta$.
\end{proof}

\begin{corollary}[Fixed joint-grid factorization]\label{occ:fc}
For every fixed $a,b\ge1$,
\begin{equation}
 \sup_\xi\sum_{\substack{u\bmod2^a\\w\text{ odd}\bmod2^b}}
 \left|\mu_{m,\xi}(J\equiv u,D\equiv w)
       -2^{1-b}\mu_{m,\xi}(J\equiv u)\right|
 \longrightarrow0.
 \label{occ:FC}
\end{equation}
This supplies the fixed-cylinder condition for the actual killed
expected occupations at each ordinary departure.
\end{corollary}
\begin{proof}
Each cell $C_{u,w}$ and level cell $L_u$ is clopen.  For product Haar,
$\pi(C_{u,w})=2^{1-b}\pi(L_u)$.  Bound each summand by
\[
 |\mu_{m,\xi}(C_{u,w})-\pi(C_{u,w})|
 +2^{1-b}|\mu_{m,\xi}(L_u)-\pi(L_u)|.
\]
Theorem~\ref{occ:uniform-convergence} makes every term tend to zero
uniformly in $\xi$, and there are finitely many cells at fixed $a,b$.
\end{proof}

\begin{remark}\label{occ:scope}
The convergence in \eqref{occ:FC} concerns the full joint grid, rather
than its separate marginals, and retains the law-dependent denominator
$a_m(\xi)$ of \eqref{occ:occupation}.  The parameters $a,b$ and the continuous input
test are fixed before $m\to\infty$.  Uniformity in the joint initial law
permits moving ordinary starting levels and offset laws; it provides neither
a rate nor a conclusion for observer precision growing with $m$.  These results are
for the canonical fresh-bit kernel and its killed expected occupation.
They do not refresh the bits of a deterministic integer trajectory or
assert deterministic Collatz no-divergence.
\end{remark}

% ---- exact/clock ----
\section{From epoch occupation to the physical excursion clock}
\label{clk:section}

Fix an odd signed integer $d$ and start the canonical pair process at
$(0,d)$. All probabilities in this section refer to its fresh fair base
bits. The positive excursion starts with the departure step and ends with
the next arrival at level zero. Its raw shortcut length $R$ includes the
departure, every positive-level holding step, and the terminal down step;
it excludes subsequent level-zero coasts. Let $K$ be the number of
positive-level direction steps. Their levels form a simple symmetric walk
from one, stopped on first hitting zero, so $K$ is odd almost surely. There
are
\begin{equation}\label{clk:counts}
 \tau_{\mathrm e}=\frac{K+1}{2}\quad\hbox{odd-level landings},\qquad
 V=\tau_{\mathrm e}-1=\frac{K-1}{2}\quad\hbox{positive even-level landings}.
\end{equation}
The odd-level, odd-offset epochs are
$Z_v=(s_v,D_v)=(2J_v+1,D_v)$, $0\le v<\tau_{\mathrm e}$. Epoch zero follows
the departure and its holding segment. Write $\mathcal G_v$ for the
information immediately before its direction bit. Add a cemetery state at
$v=\tau_{\mathrm e}$, and extend the filtration and killed process thereafter. The level
kernel is
\begin{equation}\label{clk:killed}
 \begin{split}
 j\ge1 &: \quad \PP(J_{v+1}=j-1,j,j+1\mid\mathcal G_v)
                    =\tfrac14,\tfrac12,\tfrac14,\\
 j=0 &: \quad \PP(J_{v+1}=\partial,0,1\mid\mathcal G_v)
                    =\tfrac12,\tfrac14,\tfrac14.
 \end{split}
\end{equation}
In particular $\tau_{\mathrm e}$ is a stopping time for $(\mathcal G_v)$.

We use Corollary~\ref{occ:fc} with its actual expected-occupation
normalization:
\begin{equation}\label{clk:occupation}
 a_m=\E_d(\tau_{\mathrm e}\wedge m)=2m\binom{2m}{m}4^{-m},\qquad
 \mu_m=\frac1{a_m}\E_d\sum_{v<\tau_{\mathrm e}\wedge m}\delta_{(J_v,D_v)}.
\end{equation}
Thus for every fixed $a,b\ge1$,
\begin{equation}\label{clk:fc}
 \sum_{u\bmod2^a}\ \sum_{\substack{w\bmod2^b\\w\ \text{odd}}}
 \left|\mu_m(J=u,D=w)-2^{1-b}\mu_m(J=u)\right|\longrightarrow0.
\end{equation}
The product comparison is with $\mu_m$'s own level marginal. It is not
conditioning on survival to $m$, nor averaging excursion proportions.
Only fixed precisions and the fixed departure $d$ will be used.

\begin{theorem}[The physical excursion tail]\label{clk:excursion-tail}\label{clock:tail}
For every fixed odd integer $d$, including negative $d$,
\begin{equation}\label{clk:tail-limit}
 \lim_{t\to\infty}\sqrt t\,\PP_d(R>t)=\frac2{\sqrt\pi}.
\end{equation}
More precisely, for every $q>0$,
\begin{equation}\label{clk:relative-error}
 \PP_d\bigl(|R-2(K+1)|>qK,\ K>n\bigr)=o(n^{-1/2}).
\end{equation}
No uniform convergence over $d$ is asserted.
\end{theorem}

We give the proof in several steps. The holding identity is the usual
shortcut parity correspondence \cite{Terras1976,Lagarias1985}; the walk estimates
use reflection and the central binomial formula \cite{Feller}. The
arithmetic, boundary, and occupation payments needed for this clock are
included below.

\subsection{The exact holding law and its predictable bias}

At a positive-level landing $(s,D)$, let $H$ be the number of subsequent
holding steps, before the offset next becomes odd. Put
$e_s=v_2(3^s-1)$ and take $v_2(0)=+\infty$.

\begin{lemma}[Holding law]\label{clk:holding}
Conditionally on the landing information,
\begin{equation}\label{clk:holding-law}
 H=\begin{cases}
 v_2(D),&v_2(D)<e_s,\\
 e_s+G,&v_2(D)\ge e_s,
 \end{cases}
 \qquad \PP(G=r)=2^{-r-1}\quad(r\ge0).
\end{equation}
Here $\E G=1$, $\E G^2=3$, and the conditional mean is
\begin{equation}\label{clk:eta}
 \eta(s,D)=\begin{cases}v_2(D),&v_2(D)<e_s,\\e_s+1,&v_2(D)\ge e_s.
 \end{cases}
\end{equation}
Moreover
\begin{equation}\label{clk:lte}
 e_s=1\quad(s\ \text{odd}),\qquad e_s=2+v_2(s)\quad(s\ \text{even}).
\end{equation}
If the landing offset is Haar on $\Ztwo$, then $H$ itself has the law of
$G$, independently of $s$.
\end{lemma}
\begin{proof}
The current oriented pair satisfies $Y=3^sX+D$, with $X$ Haar
conditionally on the landing information. During a holding step $X,Y$
have the same parity, and their difference is divided by two or multiplied
by three and divided by two. Its valuation therefore drops by one. Thus
\[
 H=v_2(Y-X)=v_2((3^s-1)X+D).
\]
If $v_2(D)<e_s$, this is $v_2(D)$ identically. Otherwise write
$(3^s-1)X+D=(3^s-1)(X-c)$ with
$c=-D/(3^s-1)\in\Ztwo$. The valuation of $X-c$ has the displayed
geometric law. Equal pairs form a Haar-null exception. Formula
\eqref{clk:lte} follows from $3^s\equiv3\pmod4$ at odd $s$ and
the lifting-the-exponent identity at even $s$. Under a Haar offset,
$\PP(v_2(D)=r)=2^{-r-1}$; for $r\ge e_s$ the probability of $H=r$
is $2^{-e_s}2^{-(r-e_s)-1}$, and for $r<e_s$ it is $2^{-r-1}$.
\end{proof}

\begin{lemma}[Odd landings and even predictable means]\label{clk:balance}
The holding after an odd-level landing has conditional mean one and
variance two given the past before that landing's direction bit. This
includes the initial departure. At an odd-level epoch $(s,D)$ with
$s\ge3$, define
\[
 D_+=\frac{3D+1}{2},\qquad D_- =\frac{D-3^{s-1}}2,\qquad
 \bar\eta(s,D)=\frac{\eta(s+1,D_+)+\eta(s-1,D_-)}2.
\]
Then $\bar\eta$ is the conditional mean of the next even-landing
holding. Its Haar-odd mean is one. At $s=1$ the down branch has holding
zero and
\begin{equation}\label{clk:zero-bias}
 \bar\eta(1,D)=\tfrac12\eta(2,D_+)\le2,
 \qquad |\bar\eta(1,D)-1|\le1.
\end{equation}
For $j\ge1$, put
\begin{equation}\label{clk:emax}
 E(j)=\max(e_{2j},e_{2j+2})=3+v_2(j(j+1)).
\end{equation}
Then $\bar\eta(2j+1,D)\le(2+E(j))/2$.
\end{lemma}
\begin{proof}
From an even positive level, the two odd landings have offsets
$(3D+1)/2$ and $(D-3^{s-1})/2$. Since $3^{s-1}\equiv3\pmod4$,
exactly one is even. At departure the two offsets are
$(3D+1)/2$ and $(1-3D)/2$, and again exactly one is even.
The direction bit is fair. At an odd landing $e_s=1$, so $H=0$ on
the odd branch and $H=1+G$ on the even branch. Consequently
$\E H=1$ and $\E H^2=\tfrac12(1+2+3)=3$.

For odd $s\ge3$, $3^{s-1}\equiv1\pmod8$. If $D\equiv3\pmod4$,
both $D_+,D_-$ are odd. If $D\equiv1\pmod8$, then
$v_2(D_+)=1$ and $v_2(D_-)\ge2$; if $D\equiv5\pmod8$, their
roles are reversed. Hence the mean is zero in the first case, and at
most $(1+(E(j)+1))/2$ in the second. The even-level formula
\eqref{clk:lte} gives \eqref{clk:emax}. Finally, writing $D=1+2x$
shows that both maps $D\mapsto D_+,D_-$ push Haar on the odd
integers to Haar on $\Ztwo$; Lemma~\ref{clk:holding} gives their mean
one. At $s=1$ only the up branch supplies a positive-level holding and
$e_2=3$, proving \eqref{clk:zero-bias}.
\end{proof}

For $k\ge3$, replace $e_s$ in \eqref{clk:eta} by $e_s\wedge k$ to
obtain $\eta_k$. Its Haar mean remains one. For odd $s\ge3$, set
\begin{equation}\label{clk:psi}
 \psi_k(s,D)=\frac{\eta_k(s+1,D_+)+\eta_k(s-1,D_-)}2-1,
 \qquad B_k=k+1.
\end{equation}
This has $|\psi_k|\le B_k$, Haar-odd mean zero at each $s$, and depends
only on $(s\bmod2^k,D\bmod2^{k+2})$. Also
\begin{equation}\label{clk:cap-error}
 |\bar\eta(2j+1,D)-1-\psi_k(2j+1,D)|
       \le E(j)\one_{\{E(j)>k\}},\qquad j\ge1.
\end{equation}
Indeed changing a cap $e>k$ changes \eqref{clk:eta} by at most $e$,
and changing it when $e\le k$ changes nothing.

Let $W^o$ be the sum of the odd-landing holdings minus their conditional
means, and let $W^e$ be the analogous martingale for the even landings.
In the latter, include at every $v<\tau_{\mathrm e}$ the random holding
$\widehat H_v$, which is zero when the direction step kills the excursion.
Thus $\E(\widehat H_v\mid\mathcal G_v)=\bar\eta(Z_v)$ even at the
last epoch. The exact identity is
\begin{equation}\label{clk:decomposition}
 R-2(K+1)=W^o_{\tau_{\mathrm e}}+W^e_{\tau_{\mathrm e}}+
                  \sum_{v<\tau_{\mathrm e}}(\bar\eta(Z_v)-1).
\end{equation}
There are $K+1=2\tau_{\mathrm e}$ direction steps, $\tau_{\mathrm e}$ odd holdings, and $\tau_{\mathrm e}-1$ actual
even holdings, but the last zero holding is retained in the conditional
mean and martingale decomposition. The odd martingale includes the
departure holding. Its increment variances are two. For the even
martingale they are at most $C(1+\log(v+2))^2$: the holding law is
bounded in second moment by $e^2+2e+3$, and any even landing at epoch
$v$ has level at most $2v+2$. These conditional estimates use information
available at entry, without conditioning on future excursion length.

\subsection{Level correctors and the small-excursion payment}

We record the elementary martingale estimate used repeatedly. If a
martingale, stopped at a killing time, has conditional increment second
moments at most $C(1+\log(r+2))^p$, Doob's inequality on
$[0,2^{i+1}n]$ and a union bound give
\begin{equation}\label{clk:doob}
 \PP\left(\sup_{m\ge n}\frac{|M_m|}{m}>\varepsilon\right)
 \le \frac{C_p(1+\log n)^p}{\varepsilon^2n}.
\end{equation}
Indeed the $i$th dyadic block contributes at most
$C_p(1+\log(2^{i+1}n))^p/(\varepsilon^2 2^in)$, whose sum has
the asserted order. The statement also applies when an entire holding
segment is one increment, using its entry and completion filtrations.

\begin{lemma}[Level Poisson correctors]\label{clk:level-corrector}
Suppose $f:\N\cup\{0\}\to\R$ and $\bar f\in\R$ satisfy
\[
 \left|\sum_{j=1}^x(f(j)-\bar f)\right|
       \le C(1+\log(x+1))^2\qquad(x\ge1).
\]
Then, uniformly in the departure offset, for $n\ge2$,
\begin{equation}\label{clk:level-bound}
 \PP_d\left(\tau_{\mathrm e}\ge n,\
  \left|\sum_{v<\tau_{\mathrm e}}f(J_v)-\bar f \tau_{\mathrm e}\right|>\varepsilon \tau_{\mathrm e}\right)
 \le \frac{C_f(1+\log n)^4}{\varepsilon^2n}.
\end{equation}
If the displayed partial sums are bounded independently of $x$, the
logarithmic factor may be omitted.
\end{lemma}
\begin{proof}
Replace $f$ by $f-\bar f$ and define
$g(0)=g(1)=g(\partial)=0$ and
\[
 g(j+1)-g(j)=4\sum_{i=1}^jf(i)\qquad(j\ge1).
\]
At an interior level $(Pg-g)(j)=f(j)$ by \eqref{clk:killed}. At zero
$(Pg-g)(0)=0$. With
$\xi_v=g(J_{v+1})-\E(g(J_{v+1})\mid\mathcal G_v)$, telescoping
from $J_0=0$ to the cemetery gives
\[
 \sum_{v<\tau_{\mathrm e}}f(J_v)=f(0)N_0-\sum_{v<\tau_{\mathrm e}}\xi_v,
 \qquad N_0=\#\{v<\tau_{\mathrm e}:J_v=0\}.
\]
Since $J_v\le v$, $|\xi_v|\le C_f(1+\log(v+2))^2$. Equation
\eqref{clk:doob} with $p=4$ bounds the martingale. Every visit to zero
is followed by killing with conditional probability $1/2$, so
$\PP(N_0>r)\le2^{-\lfloor r\rfloor}$. This pays the remaining
boundary term on $\tau_{\mathrm e}\ge n$. Bounded partial sums give bounded
increments and $p=0$.
\end{proof}

The useful instances have explicit arithmetic content. With $s_2(x)$
the sum of the binary digits of $x$, Legendre's formula gives
\begin{equation}\label{clk:legendre-epoch}
 \sum_{j=1}^x(E(j)-5)=1-s_2(x)-s_2(x+1).
\end{equation}
For the cap tail $f_k(j)=E(j)\one_{\{E(j)>k\}}$, $j\ge1$, take
$f_k(0)=0$. The two consecutive integers $j,j+1$ have exactly one even
member; the limiting frequency of $v_2(j(j+1))=i$ is $2^{-i}$,
$i\ge1$. Hence
\begin{equation}\label{clk:ak}
 \bar f_k=a_k=\sum_{i\ge k-2}(i+3)2^{-i}
                   =(k+2)2^{3-k}\longrightarrow0.
\end{equation}
For each $i$ the count up to $x$ differs from $x2^{-i}$ by $O(1)$.
Summing $(i+3)$ over $i\le\log_2(x+1)$ costs $O_k(\log^2(x+1))$;
the remaining infinite mean tail, multiplied by $x$, costs
$O_k(1+\log(x+1))$. This verifies the lemma's hypothesis. Finally,
$f(j)=\one_{\{j\le C\}}$ has $\bar f=0$ and bounded partial sums.
Thus cap tails are at most $(a_k+\varepsilon)\tau_{\mathrm e}$ and visits to any
fixed finite level interval are at most $\varepsilon \tau_{\mathrm e}$, apart from
events of order $O_{k,C,\varepsilon}((1+\log n)^4/n)$ on $\tau_{\mathrm e}\ge n$.

\begin{lemma}[Small excursions]\label{clk:small}
For $L\ge1$ and every odd $d$,
\begin{equation}\label{clk:small-bound}
 \E_d[R\one_{\{K\le L\}}]
       \le1+\tfrac92\E[K\wedge(L+1)]\le C\sqrt L,
\end{equation}
with an absolute $C$.
\end{lemma}
\begin{proof}
Number the landings $i=0,\ldots,K$, with levels $h_0=1$, $h_K=0$.
The departure is landing zero and the terminal holding is $H_K=0$.
Let $\mathcal L_i$ be the landing information before its holding,
and let $\tau=K\wedge(L+1)$. The event $i<\tau$ is known at that
landing. On $K\le L$,
$R=1+\sum_{i<\tau}(1+H_i)$. Nonnegativity and Lemma~\ref{clk:holding}
therefore give
\[
 \E_d[R\one_{\{K\le L\}}]
       \le1+\E\sum_{i<\tau}(e_{h_i}+2).
\]
No holding has been conditioned on $K\le L$. By \eqref{clk:lte} and
Legendre's formula,
\[
 P(x):=\sum_{h=1}^x(e_h-\tfrac52)
   =-s_2(\lfloor x/2\rfloor)-\tfrac32\one_{\{x\ \text{odd}\}}.
\]
Set $g(0)=g(1)=0$ and $g(h+1)-g(h)=2P(h)$. Then
$\tfrac12(g(h+1)+g(h-1))-g(h)=e_h-\tfrac52$ for $h\ge1$.
Since $P(h)\le0$, $g\le0$. Optional stopping for the bounded index
$\tau$ is legitimate (all visited levels are at most $L+2$), and gives
\[
 \E\sum_{i<\tau}(e_{h_i}-\tfrac52)=\E g(h_\tau)-g(1)\le0.
\]
This proves the first inequality. Reflection gives
\begin{equation}\label{clk:walk-tail}
 \PP(K>2r)=b_r:=\binom{2r}{r}4^{-r},\qquad
 \PP(K>n)\sim\sqrt{\frac2{\pi n}}.
\end{equation}
The induction $b_r\le(r+1)^{-1/2}$ follows from
$(2r+1)/(2r+2)\le\sqrt{(r+1)/(r+2)}$. Consequently
\begin{equation}\label{clk:walk-trunc}
 \PP(K>n)\le\sqrt{\frac2{n+1}},\qquad
 \E(K\wedge x)\le2\sqrt{2x}\quad(x\in\N,\ x\ge1).
\end{equation}
The last bound follows by summing the preceding tail inequality.
\end{proof}

\begin{corollary}[Uniform excursion domination]\label{clk:domination}\label{clock:domination}
For every odd $d$ and $t>0$,
\begin{equation}\label{clk:domination-bound}
 \PP_d(R>t)\le\min\{1,30/\sqrt t\},\qquad
 \E_d(R\wedge x)\le60\sqrt x\quad(x>0).
\end{equation}
\end{corollary}
\begin{proof}
For $t\ge1$, put $L=\lfloor t\rfloor$. Split according to $K>L$,
and apply \eqref{clk:walk-trunc}, \eqref{clk:small-bound}, and Markov:
\[
 \PP_d(R>t)\le\sqrt{2/t}
       +\frac{1+9\sqrt{2(L+1)}}t\le\frac{30}{\sqrt t}.
\]
For $t<1$ use the probability bound one. Integrating the tail bound
$30/\sqrt t$ on $(0,x)$ proves the second assertion.
\end{proof}

\subsection{Haar Ces\`aro bias: the stopping-time replica argument}

Let $h_{\rm odd}$ be Haar probability measure on the odd elements of
$\Ztwo$. For $k\ge3$, $L\ge1$, and odd $s\ge4L+3$, define
\begin{equation}\label{clk:cesaro}
 F_{k,L}(s,D)=\frac1{L+1}\sum_{\ell=L}^{2L}
                 \E_{(s,D)}\psi_k(Z_\ell),\qquad
 H_{k,L}(s)=\int|F_{k,L}(s,D)|\,dh_{\rm odd}(D).
\end{equation}
The level stays at least three through these epochs. These are
unconditioned future expectations from the indicated state.

\begin{lemma}[Haar Ces\`aro estimate]\label{clk:haar-bias}
Writing $r_k=\lceil(k+3)/2\rceil$ and $B_k=k+1$, uniformly for the
levels in \eqref{clk:cesaro},
\begin{equation}\label{clk:haar-bound}
 \int F_{k,L}(s,D)^2\,dh_{\rm odd}(D)
      \le\frac{2(r_k+1)B_k^2}{L+1},\qquad
 H_{k,L}(s)\le b_{k,L}:=B_k\sqrt{\frac{2(r_k+1)}{L+1}}.
\end{equation}
\end{lemma}
\begin{proof}
We first describe precisely the unread-suffix fact. Fix the base bits
of a finite positive-level raw word. Each offset row is an affine map
with an odd multiplier followed by division by two. Inductively,
specifying the offset parities through raw time $n$, including its
terminal parity, specifies one initial coset modulo $2^{n+1}$.
On it the offset at time $n$ has the form
$(uD_0+c)/2^n$ with $u$ odd. Conditional Haar measure on that coset
therefore pushes forward to Haar measure on the specified parity class.
The induction just uses that at each division the next parity determines
one additional input bit, because the multiplier is odd. This argument
also applies at a stopping time whose defining event uses only the
revealed word and terminal parity: on each value $n$ and each finite
word the unused quotient of $D_0$ is still Haar. Starting from a coset of depth $q$, after $q$ raw divisions the
offset is unrestricted Haar conditional on the consumed parities and
base bits, before its next parity is revealed. Subsequent odd epochs
have Haar-odd offset conditional on their finite words. Their centered reward \eqref{clk:psi} has mean zero.

Take two replicas starting from the same $(s,D_0)$, with
$D_0\sim h_{\rm odd}$, and using independent fair bit sequences.
Let $n_j,n'_{j'}$ be their raw finishing times at epochs $j,j'$;
let $N'(n)$ count the later epochs completed by replica two by time
$n$, including an epoch at $n$. The replica identity is
\[
 \int (\E\psi_k(Z_j))(\E\psi_k(Z_{j'}))\,dh_{\rm odd}
       =\E[\psi_k(Z_j)\psi_k(Z'_{j'})].
\]
Consider the event
\[
 A_{j,j'}=\{n_j\le n'_{j'},\ j'-N'(n_j)\ge r_k\}.
\]
At time $n=n_j$ reveal $n$, $D_0\bmod2^{n+1}$, and both length-$n$
base-bit prefixes (raw indices $0,\ldots,n-1$). This stopped information determines both paths
through $n$, the first replica's level, $N'(n)$, and $A_{j,j'}$.
Conditionally, write $D_0=\rho+2^{n+1}u$ with $u$ Haar. The offsets
are
\[
 D_n=c_1+2u_1u,\qquad D'_n=c_2+2u_2u,
\]
where $u_1,u_2$ are odd and all displayed coefficients are revealed.
The first reward reads $D_n\bmod2^{k+2}$, hence only
$u\bmod2^{k+1}$. Reveal this residue as well. The first reward is
now fixed; the second offset is Haar on a coset of depth $k+2$.
The second replica's future bits, at raw indices $n,n+1,\ldots$, are
independent of this information and of its free suffix.

On $A_{j,j'}$, its target epoch is at least $r_k$ epochs ahead.
From any current raw state, reaching the next odd-level odd epoch
takes at least one raw step, and each further such epoch requires at
least two. Thus its target occurs after at least
$1+2(r_k-1)\ge k+2$ further divisions. After those divisions the
free-suffix fact supplies a Haar offset; at the target epoch its
conditional Haar-odd mean reward is zero. The target-epoch event is
measurable from its revealed word and terminal parity. The event $A_{j,j'}$
was determined before this suffix was used, so
\[
 \E[\psi_k(Z_j)\psi_k(Z'_{j'})\one_{A_{j,j'}}]=0.
\]
The event with $n'_{j'}<n_j$ and $j-N(n'_{j'})\ge r_k$ has the
same property.

For the remaining pairs, when $n_j\le n'_{j'}$ the integer $j'$
lies between $N'(n_j)$ and $N'(n_j)+r_k-1$. There are at most
$r_k+1$ choices for each $j$. In the opposite ordering there are at
most $r_k+1$ choices for each $j'$. Pathwise, at most
$2(r_k+1)(L+1)$ pairs in $[L,2L]^2$ remain. Their rewards have
product bounded by $B_k^2$. Summing the replica identity proves the
first bound; Cauchy--Schwarz proves the second.
\end{proof}

\subsection{Fixed cylinders supply actual occupation typicality}

\begin{lemma}[Typicality with the actual source]\label{clk:typicality}
For every fixed $k\ge3$ and $L\ge1$,
\begin{equation}\label{clk:typ}
 \mu_m\!\left[\one_{\{J>2L\}}
   \bigl(|F_{k,L}(2J+1,D)|-H_{k,L}(2J+1)\bigr)\right]
                  \longrightarrow0.
\end{equation}
Consequently
\begin{equation}\label{clk:theta}
 \theta_{k,L}(d):=\limsup_{m\to\infty}
       \mu_m[\one_{\{J>2L\}}|F_{k,L}(2J+1,D)|]
         \le b_{k,L}\longrightarrow0\quad(L\to\infty).
\end{equation}
\end{lemma}
\begin{proof}
Lemmas~\ref{occ:level-haar} and \ref{occ:green} give equidistribution
of $J$ modulo each fixed power of two and the bound
$\mu_m(J\le C)\le2(C+1)/a_m\to0$ for every fixed ordinary $C$.

Fix $r$ with $2^r>4L+3$, and delete the corridor
\[
 \mathcal B_{r,L}=\{J:\operatorname{dist}(J,2^r\Z)\le2L+1\}.
\]
It has $4L+3$ distinct residues, so
$\mu_m(\mathcal B_{r,L})\to(4L+3)/2^r$.
Outside it, during the next $2L$ epoch intervals every level index
$J'$ satisfies $|J'-J|\le2L$, and neither $J'$ nor $J'+1$ is
divisible by $2^r$. Each intervening even raw level is $2J'$ or
$2J'+2$, with valuation at most $r$. Lemma~\ref{clk:holding} bounds
its arbitrary-entry mean holding by $r+3$, and bounds an odd-level
holding by two. Each interval uses two direction steps. If $n_\ell$
is the raw time of the $\ell$th later epoch, repeated conditional
expectation at these actual entries gives, for any integer or 2-adic
odd starting offset,
\begin{equation}\label{clk:raw-cap}
 \E n_{2L}\le2L(r+7),\qquad
 \PP(n_{2L}>H)\le\frac{2L(r+7)}H.
\end{equation}
The initial $J>2L$ prevents killing through these epochs. This is
a deterministic initial level restriction, not postselection on survival.

For fixed $H\in\N$, truncate the observer:
\[
 F^H_{k,L}(s,D)=\frac1{L+1}\sum_{\ell=L}^{2L}
       \E_{(s,D)}[\psi_k(Z_\ell)\one_{\{n_\ell\le H\}}],
 \qquad h^H(s)=\int|F^H_{k,L}(s,D)|\,dh_{\rm odd}(D).
\]
With $Q=H+k+3$, this is a fixed joint cylinder in
$(J\bmod2^Q,D\bmod2^Q)$ outside finitely many ordinary levels.
Indeed each raw row consumes at most one offset bit, the terminal
reward consumes at most $k+2$, and $Q$ leaves a spare bit.
The rows are
$D/2$, $(3D+1-3^s)/2$, $(3D+1)/2$, and
$(D-3^{s-1})/2$ on their respective positive-level domains.
Modulo $2^Q$, every positive power $3^e$ has period dividing
$2^Q$ in its exponent. Changing $J$ by $2^Q$ changes $s$ by
$2^{Q+1}$, preserving the affine constants at every reached row.
Taking $s>H$ keeps every $H$-step word at a positive level; any
additional levels outside the observer's domain form a finite exception.
The same finite words determine the epoch indicator and $n_\ell\le H$.
Representatives at arbitrarily large levels define the cylinder on every
level residue. Its Haar-odd average $h^H$ is periodic in that residue.

The centered cylinder $b_H=|F^H_{k,L}|-h^H$ has Haar-odd mean zero
at every level residue and $|b_H|\le B_k$. Applying \eqref{clk:fc} at
the fixed pair $(Q,Q)$ yields $\mu_m b_H\to0$ after deleting
the finite exceptions. Their cost is at most $2(C+1)B_k/a_m$.
On good corridor levels \eqref{clk:raw-cap} gives
\[
 |F_{k,L}-F^H_{k,L}|\le\frac{2B_kL(r+7)}H.
\]
The Haar average satisfies the same bound, uniformly in the offset;
the centered error is therefore at most $4B_k L(r+7)/H$. On bad
corridor levels its absolute value is at most $2B_k$. We obtain
\[
 \limsup_m\left|\mu_m[\one_{\{J>2L\}}(|F_{k,L}|-H_{k,L})]\right|
 \le\frac{4B_k L(r+7)}H+\frac{2B_k(4L+3)}{2^r}.
\]
Keep $k,L,r,H$ fixed while $m\to\infty$, next let $H\to\infty$,
and finally $r\to\infty$. This proves \eqref{clk:typ};
Lemma~\ref{clk:haar-bias} gives \eqref{clk:theta}. Neither the
precision nor the truncation horizon grows with $m$.
\end{proof}

\subsection{From occupation bias to relative control on long excursions}

The endpoint convention in \eqref{clk:occupation} matters. Put
\[
 S_L(x)=\sum_{0\le u\le\min(V,x)}
            \one_{\{J_u>2L\}}|F_{k,L}(Z_u)|.
\]
The inclusive cutoff has at most one extra summand relative to
$u<\tau_{\mathrm e}\wedge m$, of size at most $B_k$. Since
$\sum_{v=0}^{m-1}\PP(\tau_{\mathrm e}>v)=\sum_{v=0}^{m-1}b_v=2mb_m$, we have
$a_m\le2\sqrt m$, and \eqref{clk:theta} implies
\begin{equation}\label{clk:prefix-mean}
 \E_d S_L(m)\le(\theta_{k,L}(d)+\varepsilon_m)a_m+B_k,
 \qquad\varepsilon_m\longrightarrow0.
\end{equation}
For $V\in[2^in,2^{i+1}n)$, monotonicity gives
$S_L(V)\le S_L(2^{i+1}n)$. Markov and \eqref{clk:prefix-mean}
therefore give, for every $\delta>0$,
\begin{equation}\label{clk:dyadic-bias}
 \limsup_{n\to\infty}\sqrt n\,
 \PP_d(V\ge n,S_L(V)>\delta \tau_{\mathrm e})
 \le\frac{G\theta_{k,L}(d)}\delta,
 \qquad G=2\sqrt2\sum_{i\ge0}2^{-i/2}=4(1+\sqrt2).
\end{equation}
The additive $B_k$ costs $O(B_k/n)$ before multiplication by $\sqrt n$.
The errors $\varepsilon_m$ tend to zero uniformly over $m\ge n$
once $n$ exceeds their fixed onset. This does not require an onset
uniform in $d$.

The Ces\`aro average has a bounded Poisson corrector, which turns this
occupation payment into control of the actual stopped sum. Let $P$ be
the full killed epoch kernel, with an absorbing cemetery, and set
$y(s,D)=\psi_k(s,D)$ for $s\ge3$, and $y=0$ at $s=1$ and at the
cemetery. For fixed
$k,L$, define
\begin{equation}\label{clk:cesaro-corrector}
 Q_Ly=\frac1{L+1}\sum_{\ell=L}^{2L}P^\ell y,\qquad
 U_L=\frac1{L+1}\sum_{\ell=L}^{2L}\sum_{r=0}^{\ell-1}P^r y.
\end{equation}
Finite telescoping gives
\[
 (I-P)U_L=y-Q_Ly,\qquad
 \|U_L\|_\infty\le2LB_k,\qquad U_L(\partial)=0.
\]
With the stopped martingale
\[
 N_m=\sum_{v<m\wedge\tau_{\mathrm e}}
       \bigl(U_L(Z_{v+1})-PU_L(Z_v)\bigr),
\]
the exact prefix identity is
\begin{equation}\label{clk:corrector-prefix}
 \sum_{v<m\wedge\tau_{\mathrm e}}y(Z_v)
 =U_L(Z_0)-U_L(Z_{m\wedge\tau_{\mathrm e}})+N_m
                 +\sum_{v<m\wedge\tau_{\mathrm e}}Q_Ly(Z_v).
\end{equation}
The increments of $N$ are bounded by $4LB_k$. On $J_v>2L$ the next
$2L$ epochs cannot reach level zero, so $Q_Ly(Z_v)=F_{k,L}(Z_v)$
uses precisely the unconditioned future kernel in \eqref{clk:cesaro}.
At every other live state $|Q_Ly|\le B_k$. At killing the terminal
corrector is zero. Thus
\begin{equation}\label{clk:lag-boundary}
 \left|\sum_{v<\tau_{\mathrm e}}y(Z_v)\right|
 \le2LB_k+|N_{\tau_{\mathrm e}}|+S_L(V)
                  +B_k\#\{v<\tau_{\mathrm e}:J_v\le2L\}.
\end{equation}
Equation~\eqref{clk:doob} and Lemma~\ref{clk:level-corrector} pay
the martingale and the low-level count by $\varepsilon\tau_{\mathrm e}$
each, outside $O_{k,L,\varepsilon}(1/n)$ on $\tau_{\mathrm e}\ge n$.
The same identity also controls every prefix: its two corrector
endpoints cost at most $4LB_k$, independently of the prefix length.

Apply \eqref{clk:doob} to the two holding martingales in
\eqref{clk:decomposition}; apply Lemma~\ref{clk:level-corrector}
to \eqref{clk:cap-error}; and pay the zero-level bias
\eqref{clk:zero-bias} by the geometric tail of $N_0$. For fixed
$k,L,\varepsilon$, all exceptional probabilities are
$O_{k,L,\varepsilon}((1+\log n)^4/n)=o(n^{-1/2})$ on $V\ge n$.
On their complement, and when $S_L(V)\le\varepsilon \tau_{\mathrm e}$,
\begin{equation}\label{clk:clock-error}
 |R-2(K+1)|\le(7\varepsilon+a_k)\tau_{\mathrm e}+2L B_k.
\end{equation}
The seven payments are the two holding martingales, the cap-tail error,
the zero-level count, the corrector martingale, the low-level count, and
$S_L(V)$. This count leaves the limiting coefficient unambiguous.

To prove \eqref{clk:relative-error}, fix $q>0$. Choose
$0<\varepsilon<q/10$, then $k$ with $a_k\le\varepsilon$.
For each fixed $L$, when $n$ is sufficiently large the term $2L B_k$
is at most $\varepsilon \tau_{\mathrm e}$ on $V\ge n$. Since $\tau_{\mathrm e}\le K$,
\eqref{clk:clock-error} is then at most $9\varepsilon K<qK$.
Equation~\eqref{clk:dyadic-bias} bounds the remaining limiting
$\sqrt n$-probability by $Gb_{k,L}/\varepsilon$. Let $n\to\infty$
first and then $L\to\infty$. This bound vanishes. Replacing
$V\ge n$ by $K>n$, using $V=(K-1)/2$, changes only a finite
constant in the tail scale and proves \eqref{clk:relative-error}.
In particular dividing by \eqref{clk:walk-tail} gives
$R/[2(K+1)]\to1$ in probability conditionally on $K>n$; the
joint estimate \eqref{clk:relative-error} is what the tail proof needs.

\subsection{The two-sided tail squeeze}

We finish the proof of Theorem~\ref{clk:excursion-tail}. Lemma
\ref{clk:small} and Markov show that, for fixed $\eta>0$,
\begin{equation}\label{clk:small-tail}
 \limsup_{t\to\infty}\sqrt t\,
       \PP_d(R>t,K\le\eta t)\le C\sqrt\eta.
\end{equation}
Fix $0<q<1$ and $0<\eta<1/(2+q)$. On $K>\eta t$, outside the
exceptional event \eqref{clk:relative-error} with
$n=\lfloor\eta t\rfloor$,
\[
 (2-q)(K+1)\le R\le(2+q)(K+1).
\]
The error event is $o(t^{-1/2})$. Comparing the two walk thresholds
with \eqref{clk:walk-tail}, and adding \eqref{clk:small-tail} only
on the upper side, yields
\[
 \liminf_t\sqrt t\,\PP_d(R>t)\ge\sqrt{\frac{2(2-q)}\pi},
 \qquad
 \limsup_t\sqrt t\,\PP_d(R>t)
       \le\sqrt{\frac{2(2+q)}\pi}+C\sqrt\eta.
\]
Let $\eta\downarrow0$, then $q\downarrow0$, proving
\eqref{clk:tail-limit}. The proof pays the initial holding and
terminal step explicitly and uses no mixing conditional on future
survival.

% ---- exact/renewal ----
\section{Chronological intervals and their expected weight}\label{sec:assembly}

The return state, excursion length, and future departure count are
dependent. The useful quantity is the expected sum of the tail weights
at actual entries. The following elementary assembly argument applies
both to level-zero departures and to an initial return from a random
positive level.

\begin{lemma}[Weighted chronological assembly]\label{renew:weighted}
Let $M<\infty$ almost surely and let $R_1,\ldots,R_M$ be finite
nonnegative lengths, with nonnegative actual starting times $\theta_i$ and entry
histories $\mathcal F_i$. Write $B_i=\{M\ge i\}$. The event $B_i$ is
$\mathcal F_i$-measurable; on it the start $\theta_i<\infty$ and a
nonnegative finite weight $w_i$ are known, as are all earlier lengths.
Suppose
\[
 A=C+\sum_{i\le M}R_i,\qquad C\ge0,\qquad
 \theta_i+R_i\le A,\qquad
 \theta_j\ge\theta_i+R_i\quad(i<j\le M).
\]
Assume that, on each occurrence, the actual conditional tail obeys
\begin{equation}\label{renew:conditional}
 q_i(t):=\PP(R_i>t\mid\mathcal F_i),\qquad
 \sqrt t\,q_i(t)\longrightarrow c w_i\quad\hbox{almost surely},
 \qquad q_i(t)\le\frac{C_0w_i}{\sqrt t}\quad(t>0),
\end{equation}
where $c,C_0>0$ are constants independent of $i$, and the conditional
tails have versions satisfying these bounds for all $t>0$. Finally suppose
\begin{equation}\label{renew:budget}
 \E W<\infty,\qquad W=\sum_{i\le M}w_i,\qquad
 \sqrt t\,\PP(C>t)\longrightarrow0.
\end{equation}
Then
\begin{equation}\label{renew:tail}
 \sqrt t\,\PP(A>t)\longrightarrow c\E W.
\end{equation}
An integrable $C$ satisfies the remainder hypothesis. More generally it
allows an integrable coast together with a prefix having a moment of
order greater than $1/2$.
\end{lemma}

In applications \eqref{renew:conditional} follows from a fixed-state
kernel: $q_i(t)=\phi_t(X_i)$ at the actual entry, where $X_i$ is a finite
ordinary state, and $\sqrt t\,\phi_t(x)\to c w(x)$ for every fixed $x$.
Each realized $X_i$ is held fixed as $t\to\infty$; the envelope is uniform
in $x$. Thus pointwise conditional convergence holds even for an
unbounded random entry state. No uniform onset over entry states, or
independence of $M$ from the lengths, is required. A marginal tail for
$R_i$ would not supply \eqref{renew:conditional} after an earlier
interval has been observed.

\begin{proof}
Conditional tail integration gives the truncated-moment payment
\begin{equation}\label{renew:truncation}
 \E(R_i\wedge x\mid\mathcal F_i)
       =\int_0^xq_i(u)\,du\le2C_0w_i\sqrt x
                 \quad\hbox{on }B_i.
\end{equation}
The chronological long-interval events
\begin{equation}\label{renew:events}
 E_i(t)=\{M\ge i,\ \theta_i\le t,\ R_i>t\}
\end{equation}
are disjoint: an interval longer than $t$ that starts by $t$ forces every
later start past $t$. Each event implies $A>t$. Since the start and
occurrence are known at entry,
\[
 \sqrt t\sum_i\PP(E_i(t))=
 \sum_i\E[\one_{\{M\ge i,\theta_i\le t\}}\sqrt t\,q_i(t)]
              \longrightarrow c\E W.
\]
For each realized occurrence the start is finite and the conditional
limit is \eqref{renew:conditional}. The integrable dominator on counting
measure times probability is $C_0\one_{B_i}w_i$, whose integral is
$C_0\E W$. This proves the lower bound.

For the upper bound, fix $\varepsilon,\rho\in(0,1)$, put
$z=(1-\rho)t$, and take $\delta=\varepsilon^4$. Apart from
$\{C>\rho t\}$, whose probability is $o(t^{-1/2})$, the event $A>t$
implies $\sum R_i>z$. This is covered by
\begin{align*}
 E^{(1)}&=\{\exists i\le M:R_i>(1-\varepsilon)z\},\\
 E^{(2)}&=\{\exists i<j\le M:R_i>\delta z,\ R_j>\delta z\},\\
 E^{(3)}&=\left\{\sum_{i\le M}R_i\one_{\{R_i\le\delta z\}}
                                              >\varepsilon z\right\}.
\end{align*}
Outside the first two events at most one term exceeds $\delta z$, and
it is at most $(1-\varepsilon)z$, which proves this cover.
A union bound and the same weighted dominated convergence give
\begin{equation}\label{eq:assembly-large}
 \limsup_t\sqrt t\,\PP(E^{(1)})
 \le\frac{c\E W}{\sqrt{(1-\varepsilon)(1-\rho)}}.
\end{equation}

For the second event, charge each later index $j$ once. The event
$L_{<j}(x)=\{\exists i<j:R_i>x\}$, restricted to $B_j$, is known at
its entry. Consequently
\begin{equation}\label{renew:one-charge}
 \PP(E^{(2)})\le\frac{C_0}{\sqrt{\delta z}}
    \E\sum_{j\le M}w_j\one_{L_{<j}(\delta z)}
 \le\frac{C_0}{\sqrt{\delta z}}
    \E[W\one_{\{\max_{i\le M}R_i>\delta z\}}]
                         =o(t^{-1/2}).
\end{equation}
The realized collection has finite lengths, so its maximum is finite
(the empty maximum is zero). Dominated convergence uses only
$\E W<\infty$. Neither the number of earlier indices nor a product of
weights enters this payment.

Finally, \eqref{renew:truncation}, conditional expectation at each entry,
and Tonelli give
\[
 \E\sum_{i\le M}R_i\one_{\{R_i\le\delta z\}}
                 \le2C_0\E W\sqrt{\delta z}.
\]
Markov's inequality therefore yields
\begin{equation}\label{eq:assembly-small}
 \limsup_t\sqrt t\,\PP(E^{(3)})
 \le\frac{2C_0\E W\sqrt\delta}{\varepsilon\sqrt{1-\rho}}
 =\frac{2C_0\E W\varepsilon}{\sqrt{1-\rho}}.
\end{equation}
Let $t\to\infty$ first and then $\varepsilon,\rho\downarrow0$.
The upper bound agrees with the lower bound, proving \eqref{renew:tail}.
For an integrable remainder Markov gives the required negligible tail;
for a prefix moment of order $p>1/2$ it gives $O(t^{-p})$.
\end{proof}

\begin{corollary}[Predictable marks on the long interval]\label{renew:marks}
If $c\E W>0$, let $J_t$ be the unique index whose event $E_i(t)$ occurs,
assigning any value when none occurs. Then
\begin{equation}\label{renew:selection}
 \PP\left(\bigcup_iE_i(t)\mid A>t\right)\longrightarrow1.
\end{equation}
For any uniformly bounded real marks $f_i$ measurable at their entries,
\begin{equation}\label{renew:mark-law}
 \E(f_{J_t}\mid A>t)\longrightarrow
                 \frac{\E\sum_{i\le M}w_if_i}{\E W}.
\end{equation}
Off the union in \eqref{renew:selection}, the selected mark can be any
uniformly bounded value.
\end{corollary}
\begin{proof}
The proof of the lemma gives the union mass the same asymptotic as
$\PP(A>t)$, proving \eqref{renew:selection}. Conditional expectation at
entry and weighted dominated convergence, now with dominator
$C_0\sup_i\|f_i\|_\infty\one_{B_i}w_i$, give
\[
 \sqrt t\,\E\sum_i f_i\one_{E_i(t)}
                   \longrightarrow c\E\sum_{i\le M}w_if_i.
\]
The omitted part has negligible conditional probability.
\end{proof}

\subsection{The canonical departure count}

For departure $i$, let $d_i$ be its odd integer offset, $\theta_i$ its
physical starting time, and $\mathcal F_i$ the history immediately
before its departure step. Occurrence, offset, and start are known in
this history, and future bits are fresh. Put
$\phi_x(d)=\PP_d(R>x)$, $c=2/\sqrt\pi$, and $C_0=30$.
Theorem~\ref{clk:excursion-tail} and Corollary~\ref{clk:domination} give
\begin{equation}\label{eq:assembly-input}
 \sqrt x\,\phi_x(d)\longrightarrow c\quad\hbox{for every fixed odd }d,
 \qquad \phi_x(d)\le\min(1,30/\sqrt x),
\end{equation}
\begin{equation}\label{eq:truncated-excursion}
 \E_d(R\wedge x)\le60\sqrt x,
\end{equation}
while the stopping-time fresh-bit law is exactly
\begin{equation}\label{eq:departure-conditioning}
 \PP(R_i>x\mid\mathcal F_i)=\phi_x(d_i)
                       \quad\hbox{on }\{M\ge i\}.
\end{equation}
The count here includes every actual departure, including those inside
compact absorbing-word trials. The earlier reward estimates give
\begin{equation}\label{eq:assembly-moments}
 1\le\E M<\infty,\qquad \E C_{\mathrm{coast}}<\infty,\qquad
 A=\sum_{i\le M}R_i+C_{\mathrm{coast}}.
\end{equation}
The separately established second departure moment remains available
for the triple-sector argument; it is not a scalar assembly hypothesis.

\begin{proof}[Proof of the shortcut assertion in Theorem~\ref{thm:exact}]
Apply Lemma~\ref{renew:weighted} with $w_i=1$ at each actual departure
and $C=C_{\mathrm{coast}}$. Equations
\eqref{eq:assembly-input}--\eqref{eq:assembly-moments} supply precisely
its conditional laws and budgets. Hence
\begin{equation}\label{eq:assembly-lower}
 \sqrt t\,\PP(A>t)\longrightarrow\kappa=c\E M\in(0,\infty).
\end{equation}
\end{proof}

\subsection{Initial returns and random entrances}

Suppose the process first enters at a finite stopping time $\Theta$
with an ordinary random state of positive level $S$, and the initial
return $R_0$ has conditional coefficient $cS$ and envelope
$C_1S/\sqrt t$. After that return its $M_*$ departure excursions have
coefficient $c$ and envelope $30/\sqrt t$. Include a zero initial
interval when $S=0$. If
\[
 \E S+\E M_*+\E C_*<\infty,\qquad
 \sqrt t\,\PP(\Theta>t)\longrightarrow0,\qquad
 A=\Theta+R_0+C_*+\sum_{i\le M_*}R_i,
\]
Lemma~\ref{renew:weighted}, after relabelling interval zero as the first
index, gives
\begin{equation}\label{renew:initial-return}
 \sqrt t\,\PP(A>t)\longrightarrow c\E(S+M_*).
\end{equation}
Use weights $S,1,\ldots,1$, envelope constant $\max(C_1,30)$, and
remainder $\Theta+C_*$. Its tail is negligible by a union bound.
The initial state's law is the actual entrance law: conditioning there
holds each finite realized state fixed, and its linear envelope is
integrable because $\E S<\infty$. In particular no moment of $S^2$,
$SM_*$, or $M_*^2$ is required. This single specialization supplies the
scalar assembly both for a fixed positive starting level and for the
random two-class entrance of a finite family.

% ---- exact/capacity ----
\section{Scalar capacity, conditioned histories, and the absorption domain}
\label{harm:section}

The renewal theorem turns one scalar observable into the coefficient of
survival. From an integer canonical state $(s,D)$ let $M_{s,D}$ count every future
level-zero odd departure, including those in compact trials, and set
\begin{equation}\label{harm:capacity}
 h_{\rm cap}(s,D)=s+\E_{s,D}M_{s,D},\qquad h_{\rm cap}(0,0)=0.
\end{equation}
For $z=(s,D)$ abbreviate $M_z=M_{s,D}$ and $h_{\rm cap}(z)=h_{\rm cap}(s,D)$.
The initial level weighs the first positive-level return; each subsequent
departure has weight one. This notation leaves the ordinary halving
count $H(x,m)$ untouched. We first supply the scalar return law needed
for this interpretation, independently of any family theorem.

\subsection{The scalar positive-level supplier}

\begin{lemma}[Return from a fixed positive level]\label{pair:fixed-return}\label{family:fixed-return}
Start from any fixed canonical state $(s,D)$ with $s\ge1$, integer
$D$, and a fresh Haar base. If $R_{s,D}$ is its first physical return to level zero, then
\[
 \sqrt t\,\mathbb P(R_{s,D}>t)\longrightarrow cs,
 \qquad c=2/\sqrt\pi,
 \qquad \mathbb P(R_{s,D}>t)\le\min(1,Cs/\sqrt t).
\]
No onset uniform in an unbounded initial level is asserted.
\end{lemma}
\begin{proof}
The domination is Lemma~\ref{pair:return-bound}. We give the
adapter needed for the exact coefficient, retaining
the source occupation normalization.

For odd $s=2j+1$ and odd offset, the killed epoch chain of
\eqref{occ:killed-level} starts at the fixed $j$. Its Green function is
\[
 G_j(k)=2\min(2j+1,2k+1).
\]
This follows by multiplying its affine hitting probability by the mean
$2(2k+1)$ visits from $k$. In particular its mean zero visits is still
two. If $K_s$ is the direction count to zero, reflection gives
\[
 \mathbb P(K_s>n)\sim s\sqrt{2/(\pi n)},\qquad
 \mathbb P(K_s>n)\le\min(1,Cs/\sqrt{n+1}).
\]
The epoch lifetime is $(K_s+1)/2$, and its expected prefix occupation
normalizer therefore satisfies
\[
 a_m(j)\sim2s\sqrt{m/\pi},\qquad a_m(j)\le Cs\sqrt m.
\]
The periodic level corrector has bounded starting value at $j$ and the
same two-visit zero boundary payment. Fixed ordinary levels have
occupation at most $2s(C+1)/a_m(j)$. Corollary~\ref{occ:fc} applies
with the actual initial law $\xi=\delta_j\otimes\delta_D$ and its
own denominator $a_m(j)$; it supplies the same fixed-grid product Haar
factorization. No denominator from the $j=0$ chain is substituted.

In the unbounded level corrector the additional starting term is the
fixed number $-g(j)$, and $J_v\le j+v$ keeps the same logarithmic
martingale variance estimates with fixed-$j$ constants. The small
excursion Legendre corrector starts at $s$ and gives
\[
 \mathbb E[R_{s,D}\mathbf1_{\{K_s\le L\}}]
 \le\tfrac92\mathbb E[K_s\wedge(L+1)]-g(s)
 \le C_s\sqrt L.
\]
The optional direction index is bounded and $-g(s)$ is finite. The
Haar-replica estimate concerns high current levels and is unchanged;
the actual-source typicality and dyadic prefix argument now use the
factorization above and $a_m(j)\le Cs\sqrt m$. The dyadic-bias
constant becomes $CsG$, still multiplied by $\theta_{k,L}\to0$;
the level-corrector payment retains $p=4$ when $J_0=j$.

There is no initial departure or initial odd-level holding. To fix the
clock convention, number the direction landings $0,1,\ldots,K_s$ and
let $H_q$ be the holding after landing $q$. At this odd-offset entrance
$H_0=H_{K_s}=0$, and the physical direction clocks are
\[
 \sigma_0=0,\qquad \sigma_q=q+\sum_{r<q}H_r,\qquad
 R_{s,D}=\sigma_{K_s}.
\]
For $q=2m$ the noninitial holdings consist of $m$ even-level holdings
and $m-1$ odd-level holdings. The odd holdings therefore number one
less than the epoch count. With the holding martingales defined as in
\eqref{clk:decomposition}, this gives the exact terminal identity
\[
 R_{s,D}-2(K_s+1)
 =W^o+W^e+\sum_{v<\tau_e}(\bar\eta(Z_v)-1)-2.
\]
The holding martingale, cap, corrector, and lagged bias payments of
Section~\ref{clk:section}, with the fixed starting terms just specified,
prove
\[
 \mathbb P(|R_{s,D}-2K_s|>qK_s,K_s>n)=o(n^{-1/2}).
\]
The small-excursion bound and the same two-sided squeeze give coefficient
$cs$ from the reflected $K_s$ tail.

An initial even offset first has a holding with mean at most $e_s+1$.
At its actual completion the offset is an odd integer and the future
base is Haar. The coefficient limits for this selected law are dominated
by $Cs/\sqrt t$, and the holding has negligible square-root tail.
For even $s\ge2$, after that holding take one direction and complete its
odd-level holding, of conditional mean at most two in either branch.
The odd entrance levels are
$s-1,s+1$ with fair probabilities. Their dominated coefficients average
to $cs$, while the initial step and holding are negligible. All entrance
mixtures use conditional laws at their actual stopping times. Explicitly,
if $L$ is one of these finite initial prefixes and $Q$ its remaining
return time, dominated convergence at the actual entrance gives the
coefficient for $Q$. For $0<\eta<1$,
\[
 \{L+Q>t\}\mathbin{\triangle}\{Q>t\}
 \subset\{L>\eta t\}\cup\{(1-\eta)t<Q\le t\}.
\]
The first event has probability $O_s(t^{-1})$, and the second has
$\sqrt t$-normalized limiting mass bounded by
$cs((1-\eta)^{-1/2}-1)$ after averaging the entrance levels.
First let $t\to\infty$ and then $\eta\downarrow0$. This pays the
prefix in the coefficient, without conditioning its holding on future
survival.
\end{proof}

\subsection{Harmonic cancellation and survival conditioning}

\begin{theorem}[Harmonic capacity and conditioned histories]\label{harm:law}
Write $K_{\rm pair}$ for the six-row kernel, including its boundary
orientation swap. The capacity is finite at every integer state and
positive at every nonabsorbed state,
and
\[
 \sqrt t\,\PP_z(A_z>t)\longrightarrow\frac2{\sqrt\pi}h_{\rm cap}(z),
 \qquad K_{\rm pair}h_{\rm cap}=h_{\rm cap}.
\]
For $z\ne(0,0)$ and every fixed integer history length $h\ge0$, its labelled-history law conditioned
on $A_z>t$ converges in total variation to
\[
 \PP_z^{h_{\rm cap}}(E)=\frac{\E_z[\one_E h_{\rm cap}(Z_h)]}{h_{\rm cap}(z)}.
\]
These consistent laws have the Doob transition kernel
$K_{\rm pair}^{h_{\rm cap}}(z,z')=K_{\rm pair}(z,z')h_{\rm cap}(z')/h_{\rm cap}(z)$ on nonabsorbed states.
\end{theorem}
\begin{proof}
At positive level, first return to zero has raw offset $U$ satisfying
\[
 \E_z|U|^\alpha\le r^s|D/3^s|^\alpha+b_\infty.
\]
The boundary reward theorem at that return gives finite future departure
mean and coast mean. At level zero these are the existing reward bounds.
Lemma~\ref{pair:fixed-return} supplies the initial interval coefficient
$2s/\sqrt\pi$ and linear domination. Subsequent departing intervals have
coefficient $2/\sqrt\pi$ and the universal envelope.
The initial-return specialization \eqref{renew:initial-return} applies with the initial return weighted
by $s$ (omitted at $s=0$), every subsequent departure weighted by one,
and integrable boundary coasts. It yields the coefficient
$2(s+\E_zM_z)/\sqrt\pi$.
No product of the initial level and future count is used. Positivity
follows from $s>0$, or from the first genuine departure after the finite
coast of a nonzero level-zero offset.

Put $b(z)=\one_{\{s=0,D\text{ odd}\}}$. The exact kernel gives
$\E_z s_1=s+b(z)$, while the full count satisfies
$\E_zM_z=b(z)+\E_z\E_{Z_1}M$. Thus the level gains exactly the
amount lost from the future departure count. All expectations are
finite, so adding these identities cancels the drift and proves
$K_{\rm pair}h_{\rm cap}=h_{\rm cap}$. The assigned value at the
absorbing state is zero.

There are only finitely many labelled histories of a fixed length $h$.
On each history the state $Z_h$ is a fixed integer state.
Markov conditioning gives, for $t>h$,
\[
 \PP_z(E,A_z>t)=
 \E_z[\one_E\one_{\{A_z>h\}}\PP_{Z_h}(A>t-h)].
\]
Apply the state tail on that finite set and divide by its value from $z$.
This proves the history limit atom by atom, hence in total variation.
Repeated harmonicity normalizes the limits and makes them consistent.
Their one-step likelihood ratio is exactly $h_{\rm cap}(z')/h_{\rm cap}(z)$.
The survival indicator may be omitted in the limiting formula because
$h_{\rm cap}(0,0)=0$. Duplicate successors are aggregated in the state kernel;
for example both labels from $(2,4)$ lead to $(2,2)$.
The driver becomes biased by capacity; no fresh fair-bit assertion under
the limiting survival-conditioned law is made.
\end{proof}

This is microscopic survival conditioning, distinct from the diffusive
meander law. The Doob construction is classical; the explicit harmonic
function comes from the canonical level drift and the expected full
departure count. The result has fixed-state, fixed-history scope and
claims no convergence rate, stationary law of $K_{\rm pair}^{h_{\rm cap}}$, prescribed integer
trajectory law, or no-divergence consequence.

\subsection{The exact arithmetic domain of absorption}\label{triad:section}

The integer-offset assumption admits a sharp extension. Denote by
$\Z[1/3]$ the rationals with denominator a power of three, embedded
in $\Ztwo$.

\begin{theorem}[Admissible offsets]\label{triad:domain}
Start the actual canonical fresh-bit kernel from $(s,D)$, where $s$ is
a nonnegative integer and $D\in\Ztwo$. If $D\notin\Z[1/3]$, no
finite label word reaches absorption. If $D\in\Z[1/3]$, absorption
is almost sure, $\E_{s,D}M<\infty$, and
\[
 \sqrt t\,\PP_{s,D}(A>t)
       \longrightarrow\frac2{\sqrt\pi}(s+\E_{s,D}M).
\]
The coefficient is zero at $(0,0)$ and positive elsewhere.
For a Haar pair $(x,x+d)$, every fixed nonzero $d\in\Z[1/3]$ has
the shortcut and ordinary coefficients $\kappa_d=2\E M_d/\sqrt\pi$
and $\beta_d=\sqrt{3/2}\kappa_d$. For fixed $d\notin\Z[1/3]$,
balanced absorption is impossible for every root, and the unrestricted
ordinary first meeting is infinite Haar-almost surely.
\end{theorem}
\begin{proof}
Every six-row step at an ordinary level has
\[
 D'=(aD+b)/2,\qquad a\in\{1,3,-3\},\quad b\in\Z.
\]
Thus a length-$h$ word has offset
$(\epsilon3^qD+B)/2^h$, $\epsilon\in\{-1,1\}$, $q\ge0$, $B\in\Z$.
Its vanishing forces $D\in\Z[1/3]$, proving necessity pointwise.

Now let $D$ be triadic. Let $r(D)$ be its denominator exponent in lowest
terms, zero for integers. The rows preserve $\Z[1/3]$: their parity
condition makes the numerator divisible by two. When $r(D)>0$, a row
with coefficient magnitude one preserves that exponent, and a row with
coefficient magnitude three reduces it by exactly one. At positive
level and at an even level-zero offset, one fair label has the latter
coefficient; at an odd level-zero offset both do. The stopping time
\[
 \tau=\inf\{h:D_h\in\Z\}
\]
is therefore dominated by the waiting time for $r(D)$ fair Bernoulli
successes. Choose predictably one clearing label at each noncleared state;
these indicators are fresh fair bits, while any additional clearing helps.
Consequently
\begin{equation}\label{triad:clearing}
 \E e^{u\tau}\le\left(\frac{e^u}{2-e^u}\right)^{r(D)}<\infty,
 \qquad 0<u<\log2.
\end{equation}

Use the real rational embedding before clearing and put
$w_h=D_h/3^{s_h}$. The normalized rows give
$|w_{h+1}|\le(3/2)|w_h|+1/2$, hence
\[
 |w_h|\le(3/2)^h(|w_0|+1).
\]
For any fixed $0<\alpha<1$, $\alpha\log(3/2)<\log2$, so
\eqref{triad:clearing} gives
\[
 \E(1+|w_\tau|^\alpha)<\infty,\qquad
 \E s_\tau\le s+\E\tau<\infty.
\]
At clearing the offset is integer. If $S=s_\tau>0$, its first raw
return offset $V$ satisfies
\[
 \E(|V|^\alpha\mid\mathcal F_\tau)
       \le r_\alpha^S|w_\tau|^\alpha+b_\infty,
 \qquad r_\alpha<1,
\]
by Lemma~\ref{pair:returnmoments}; if $S=0$, take $V=D_\tau=w_\tau$.
The integer-boundary reward theorem then gives finite future departure
mean and finite future coast mean. Integer-state absorption after this
entry is almost sure. All entry values are selected before fresh future
bits, so these conditional source estimates apply.

Let $M_{\rm pre}$ count departures starting before $\tau$, and
$M_{\rm post}$ the subsequent departures. An excursion straddling
$\tau$ is not recounted. With $R_0$ the initial positive-level return
from $S$ (zero for $S=0$), the future decomposition is
\[
 A=\tau+R_0+C_{\rm post}+\sum_{i\le M_{\rm post}}R_i.
\]
The initial-return specialization \eqref{renew:initial-return} applies with integrable $S$, finite
future count and coast means, and negligible prefix $\tau$.
It gives coefficient $c\E(S+M_{\rm post})$, $c=2/\sqrt\pi$;
the initial return uses domination $CS/\sqrt t$, and every later
charged interval is a standard integer departure. No product moment
or second count moment is required.

The exact level-minus-departure martingale, stopped at $\tau\wedge n$,
gives $\E S=s+\E M_{\rm pre}$ after dominated convergence:
its absolute stopped value is at most $s+2\tau$ and $\E\tau<\infty$.
Since $M=M_{\rm pre}+M_{\rm post}$ and $M_{\rm pre}\le\tau$,
the asserted coefficient and finite count mean follow. Positivity is
immediate at $s>0$. At level zero a nonzero offset has a finite even
coast followed by a genuine departure.

For initial level zero, the first-clock identity of
Lemma~\ref{ordinary:exactclock} and original-root parity fairness in
\eqref{ordinary:parities} use the literal kernel and Haar translation
for every fixed two-adic shift. The stopped concentration and ordinary
squeeze of Theorem~\ref{ordinary:tail} therefore transfer the triadic
shortcut coefficient to both ordinary clocks.
Outside the domain no balanced witness exists. At a fixed ordinary time $j$,
an unbalanced meeting on a fixed pair of realizing words has unequal
affine slopes $3^j/6^{H_1}\ne3^j/6^{H_2}$, so it fixes at most one two-adic root. There are finitely
many words at each time and countably many times; their union is Haar-null.
\end{proof}

The harmonic-prefix proof extends to triadic offsets. Away from absorption,
the function $h_{\rm cap}(s,D)=s+\E_{s,D}M$ is finite and its drift cancellation
remains valid. Every fixed history has only finitely many triadic
successors, so the proof of Theorem~\ref{harm:law} applies.

The earlier raw-kernel analysis already identifies the carrier
$\Z\times\Z[1/3]$. Here the scalar capacity determines the exact
absorption domain and its physical coefficient. The domain includes $d=1/3$ and excludes
$d=1/5$. The result keeps each offset fixed and gives no empirical
Brownian law for positive rational starting values.

% ---- exact/ordinary ----
\section{The first ordinary meeting clock}\label{sec:ordinary}

For the original adjacent roots $a,b=a+1$, use the ordinary clocks
$B$ and $U$ of Section~\ref{sec:intro}, including for deterministic
integer realizations.

\begin{lemma}[Exact first balanced ordinary clock]\label{ordinary:exactclock}
For distinct integer or $2$-adic roots $a,b=a+d$ with $d\ne0$, $B<\infty$ if and only
if $A<\infty$. On this event,
\[
 B=A+o_A(a)-1=A+o_A(b)-1,
 \qquad H(a,B)=H(b,B)=A-1,
 \qquad A\le B\le2A-1.
\]
On this event the first unrestricted ordinary meeting satisfies $U=B$.
Writing $\tau_h(x)=h+o_h(x)$ for the completed-grid time, one also has
\[
 A\le h\Longrightarrow B\le\tau_h(a)-1,
 \qquad h<A<\infty\Longrightarrow B\ge\tau_h(a).
\]
\end{lemma}
\begin{proof}
Inspection of the six canonical rows shows that a first transition into
$(0,0)$ is necessarily the down-step from $(1,1)$ with odd oriented base
$X$. Set $h=A-1$. Immediately before this transition the other entry is
$Y=3X+1$, which is even. The signed odd-count identity through swaps says
that the root represented by $Y$ has one more odd shortcut branch than
the root represented by $X$. Write their counts at height $h$ as $q+1$
and $q$. Their completed ordinary grids end at times $h+q+1$ and $h+q$.
At time $h+q+1$ the latter trajectory has taken its next ordinary odd
branch and is at $3X+1=Y$, while the former is at its completed endpoint
$Y$. Each has $h$ halvings and $q+1$ odd branches. Completing the next
halving gives shortcut absorption at height $A$, with common odd count
$o_A(a)=o_A(b)=q+1$. Thus a balanced merger occurs at $A+o_A(a)-1$.

Conversely, suppose a balanced ordinary merger exists and take its first
time $B$. Put $H=H(a,B)=H(b,B)$ and $R=B-H$. Its final branches cannot
both be even or both be odd: injectivity on either parity class would
give equal penultimate endpoints, and deleting equal branch types would
preserve equal counts, contradicting firstness. Thus the final branches
are mixed. The odd-final path has odd penultimate value $u$ and common
endpoint $y=3u+1$. At shortcut height $H$ this path has endpoint $u$ and
odd count $R-1$, while the even-final path has endpoint $y$ and odd count
$R$. Since $y$ is even, one further shortcut step gives equal endpoints
$y/2$ and equal odd counts $R$. The dictionary implies $A\le q:=H+1$.
If $A<q$, monotonicity of original odd counts and the merger constructed
above give a balanced merger at
\[
 A+o_A(a)-1\le q+o_q(a)-2=B-1,
\]
a contradiction. Hence $A=q$ and $B=A+o_A(a)-1$, with $H=A-1$.
Finally $1\le o_A(a)\le A$. Nothing in this argument uses positivity
or an ordering on the roots.

After any ordinary meeting, endpoints and branch types coincide forever,
so the difference of the two halving counts remains constant. Since it
is zero at $B$, it was already zero at $U\le B$, proving $U=B$.
Finally $\tau_h$ is strictly increasing with increments one or two.
Thus $B=\tau_A(a)-1$ gives both displayed comparisons. If $A=\infty$,
the equivalence already proved gives $B=\infty$.
\end{proof}

\begin{lemma}[Original-root parities under the full revealed history]
\label{ordinary:parities}
Under Haar initial law, each original trajectory's shortcut odd indicators
are iid Bernoulli$(1/2)$. Let $\mathcal F_h$ be the sigma-field generated by the canonical
labels through $h-1$, together with the symbolic states, orientations,
swaps, counts, and absorption flags determined by those labels.
Then each original root's next odd indicator is conditionally fair
 given $\mathcal F_h$.
\end{lemma}
\begin{proof}
The current oriented-base parity $\epsilon_h$ is fresh and fair given
$\mathcal F_h$. If a fixed original root is currently the base, its parity
is $\epsilon_h$. If it is the second entry, its parity is
$\epsilon_h\mathbin{\mathrm{xor}}(D_h\bmod2)$, since $3^{s_h}$ is odd.
Which case applies and $D_h\bmod2$ are $\mathcal F_h$-measurable.
Thus the original-root parity is a predictable permutation of the next
fair bit. Conditional fairness gives the iid claim by induction. Retaining
the common base after absorption proves the assertion at all heights.
The two original parity sequences are not claimed independent of each
other, and the entire current $2$-adic base is not included in the
revealed-history sigma-field.
\end{proof}

\begin{lemma}[Joint stopped-time concentration]\label{ordinary:concentration}
For $0<\delta<1/2$ and integer $n\ge0$,
\[
 \PP\bigl(A>n,\ |o_A(a)-A/2|>\delta A\bigr)
 \le \frac{2e^{-2\delta^2(n+1)}}{1-e^{-2\delta^2}}.
\]
The same bound applies to the event that the deviation occurs at any
integer height $h>n$. In particular, $\delta=n^{-1/4}$ (for $n>16$) makes the bound
$O(\sqrt n\,e^{-2\sqrt n})=o(n^{-1/2})$ as $n\to\infty$.
\end{lemma}
\begin{proof}
Hoeffding's inequality and Lemma~\ref{ordinary:parities} give
$\PP(|o_h(a)-h/2|>\delta h)\le2e^{-2\delta^2h}$.
Sum this over $h>n$. Almost-sure finiteness of $A$ makes the stopped-time
event a subset of this union. No independence from $A$ is used.
\end{proof}

\begin{theorem}[Exact ordinary meeting tail]\label{ordinary:tail}
Under Haar initial law,
\[
 \sqrt t\,\PP(B>t)\longrightarrow\sqrt{\frac32}\,\kappa,
 \qquad
 \sqrt t\,\PP(U>t)\longrightarrow\sqrt{\frac32}\,\kappa.
\]
In fact $U=B$ on $\{A<\infty\}$, and hence Haar-almost surely.
\end{theorem}
\begin{proof}
Set $\delta_t=t^{-1/4}$, and let $E_t$ be the event that
$|o_h(a)-h/2|>\delta_t h$ at some integer $h>\lfloor t/2\rfloor$.
The preceding geometric bound gives
$\PP(E_t)=O(\sqrt t\,e^{-\sqrt t})=o(t^{-1/2})$.
The bound $B\le2A-1$ implies that $B>t$ forces $A>\lfloor t/2\rfloor$.
Outside $E_t$, at such an absorption height,
\[
 (3/2-\delta_t)A-1\le B\le(3/2+\delta_t)A-1.
\]
Therefore, for all sufficiently large $t$,
\begin{align*}
 \PP\left(A>\frac{t+1}{3/2-\delta_t}\right)-\PP(E_t)
 &\le\PP(B>t)\\
 &\le\PP\left(A>\frac{t+1}{3/2+\delta_t}\right)+\PP(E_t).
\end{align*}
Both thresholds divided by $t$ tend to $2/3$. The established asymptotic
$\sqrt h\,\PP(A>h)\to\kappa$, applied to their integer parts, proves
the first assertion.

Lemma~\ref{ordinary:exactclock} gives $U=B$ on finite absorption.
Almost-sure Haar absorption therefore proves the second tail assertion.
\end{proof}

\begin{remark}[The exceptional roots and the positive-root transfer]
The equality $U=B$ has an almost-sure, rather than unrestricted pointwise,
scope. The adjacent $2$-adic roots $a=-8/5$, $b=-3/5$ satisfy
$C(a)=a/2=-4/5=3b+1=C(b)$ at time one with unequal halving counts.
Hence $U=1$ but $B=A=\infty$.

For deterministic integer roots, $\{B\le k\}$ is decided by the first
$k$ canonical labels because $A\le B$ and $o_A$ is revealed at absorption.
It therefore has exact cylinders modulo $2^k$. For positive roots,
$C^j(x)\ge x/2^{H(x,j)}$; since the first balanced merger has $A-1$
halvings, $x>2^{A-1}$ ensures that its entire prefix stays strictly above
one. These facts preserve the before-one clause needed for Gao's predicate.
\end{remark}
\begin{corollary}[Every fixed integer offset]\label{ordinary:fixedoffset}
For each fixed nonzero $d\in\Z$, start the Haar pair from $(x,x+d)$,
and let $A_d,B_d,U_d,M_d$ denote its corresponding clocks and departure
count. Then
\[
 1\le\E M_d<\infty,\qquad
 \kappa_d:=\frac{2\E M_d}{\sqrt\pi}.
\]
\begin{align*}
 \sqrt t\,\PP(A_d>t)&\longrightarrow\kappa_d,\\
 \sqrt t\,\PP(B_d>t)&\longrightarrow\sqrt{3/2}\,\kappa_d,\\
 \sqrt t\,\PP(U_d>t)&\longrightarrow\sqrt{3/2}\,\kappa_d.
\end{align*}
The equality of the ordinary tails is an almost-sure equality of their
clocks. No uniform asymptotic as $d$ grows with $t$ is asserted.
\end{corollary}
\begin{proof}
Lemma~\ref{pair:rewards} applies to every integer boundary offset $d$.
A nonzero initial even offset has a finite
coast of exactly $v_2(d)$ steps and cannot be absorbed during it; thus
$M_d\ge1$ and the total coast still has finite mean. The decomposition
$A_d=\sum_{i\le M_d}R_i+C_{\mathrm{coast}}$ and
Lemma~\ref{renew:weighted}, with unit weight at each departure,
now apply with this initial law. Each
realized departure remains one fixed signed odd integer and has the same
excursion coefficient and uniform domination.

The first-clock identity, predictable parity permutation, and the argument
that $U=B$ on finite absorption use only distinct roots with initial level
zero. They apply verbatim to $(x,x+d)$. The stopped parity concentration
therefore transfers $\kappa_d$ to both ordinary clocks.
\end{proof}

% ---- exact/populations ----
\section{One transfer to finite integer populations}\label{sec:population}

All population estimates below use the same counting fact. Its error
retains the size of the selected residue set.

\begin{lemma}[Periodic-set discrepancy]\label{pop:periodic}
Let $E\subseteq\Z$ be a union of $K$ classes modulo $Q\ge1$, and put
$p=K/Q$. For any interval $I$ of $N$ consecutive integers,
\[
 \left|\frac{\#(I\cap E)}N-p\right|\le\frac{Qp(1-p)}N,
 \qquad
 \frac{\#(I\cap E)}N\le(1+Q/N)p.
\]
For a complete union of periods the discrepancy is zero.
\end{lemma}
\begin{proof}
Write $N=\ell Q+r$, $0\le r<Q$. Complete periods contribute $\ell K$.
The remaining interval contains $a$ selected residues with
$\max(0,r-(Q-K))\le a\le\min(r,K)$. Hence
$a-rp\le K(1-p)$; applying the same inequality to the complementary
residues gives $rp-a\le(Q-K)p=K(1-p)$. Division by $N$ proves both bounds.
\end{proof}

\begin{proposition}[Exact absorption-cylinder discrepancy]\label{pop:discrepancy}
Let $B_t$ be the positive integer starts whose adjacent pair has not
undergone balanced shortcut absorption by height $t$. For every interval
$I$ of $N$ consecutive positive integers and every integer $t\ge0$,
\begin{equation}\label{eq:population-error}
 \left|\frac{\#(I\cap B_t)}N-p_t\right|
 \le\frac{2^tp_t(1-p_t)}N,\qquad p_t=\PP(A>t).
\end{equation}
Thus, if $t\to\infty$ and $2^t/N(t)\to0$,
\[
 \sqrt t\,\frac{\#(I(t)\cap B_t)}{N(t)}\longrightarrow\kappa.
\]
\end{proposition}
\begin{proof}
Lemma~\ref{pair:cylinders} identifies exactly $2^tp_t$ nonabsorbed
classes modulo $2^t$. Apply Lemma~\ref{pop:periodic}. Multiplication of
its error by $\sqrt t$ leaves at most
$(2^t/N)\sqrt t\,p_t\to0$, by the shortcut assertion of
Theorem~\ref{thm:exact}.
\end{proof}
For complete periods this gives the exact residue identity
\eqref{eq:duran-law}, hence the sequence asymptotic
\eqref{eq:sequence-asymptotic}, without a population-size restriction.
The same discrepancy bound holds for every fixed or specified integer
offset $d$, with $p_t$ replaced by $S_t(d)$.

\subsection{Gao's dyadic population}\label{sec:gao}

\begin{lemma}[Ordinary balanced cylinders]\label{gaoexact:cylinders}
For every integer $k\ge0$, the event $B(n)\le k$ is a union of residue
classes modulo $2^k$. Its earliest witness, the ordinary parity words
through time $k$, and their branch counts are fixed on each class. Thus,
with $Q_k=\PP(B>k)$,
\[
 \#\{0\le n<2^k:B(n)>k\}=2^kQ_k.
\]
\end{lemma}
\begin{proof}
A canonical word of length $k$ fixes both original shortcut parity words
through height $k$, by tracking orientation and the offset parity.
Expanding these words fixes both ordinary words through time $k$.
If $A>k$, then $B>k$ by Lemma~\ref{ordinary:exactclock}; otherwise the
word fixes $A$ and $o_A$, hence $B=A+o_A-1$ and all branch counts.
Lemma~\ref{pair:cylinders} gives one such word per residue modulo $2^k$
and Haar mass $2^{-k}$ per class. Endpoint values themselves may vary
between lifts; the witness and its counts do not.
\end{proof}

\begin{lemma}[Exponentially sparse visits to one]\label{gaoexact:one}
Let $F_0=0$, $F_1=1$ and $F_{j+2}=F_{j+1}+F_j$. For every integer $k\ge0$,
\[
 \#\{x\ge1:\exists j\le k,\ C^j(x)=1\}\le F_{k+3}-1.
\]
\end{lemma}
\begin{proof}
An ordinary odd branch is followed by an even branch when the next branch
is present. Words of length $j$ without consecutive odd branches number
$F_{j+2}$, by their final symbols. A first visit to one at time $j\ge1$
has final even branch, because two is its unique positive preimage. Its
preceding word therefore has at most $F_{j+1}$ possibilities. The affine
map for a fixed word has nonzero slope $3^R/2^H$, so at most one positive
preimage of one. Including the start one and summing over first hitting
times gives $1+\sum_{j=1}^kF_{j+1}=F_{k+3}-1$.
\end{proof}

\begin{proposition}[Exact Gao adapter]\label{gaoexact:adapter}
For Gao's strict interval $1\le n<2^k-1$, denominator $2^k$, and balanced
merger by $k$ before either trajectory visits one through that merger,
\[
 0\le(1-\bar d_k)-Q_k\le 2^{1-k}F_{k+3}
 =O((\varphi/2)^k),\qquad \varphi=(1+\sqrt5)/2.
\]
Consequently $\sqrt k(1-\bar d_k)\to\sqrt{3/2}\,\kappa$.
\end{proposition}
\begin{proof}
The complete period contains $2^k(1-Q_k)$ balanced successes. Gao removes
the two starts $0,2^k-1$. Any further lost success has one of its roots
visit one by time $k$. Lemma~\ref{gaoexact:one} bounds these losses by
$2(F_{k+3}-1)$, so the total loss is at most $2F_{k+3}$. Divide by $2^k$.
The resulting exponentially small error and Theorem~\ref{ordinary:tail}
give the limit.
\end{proof}

\subsection{All integer cutoffs and the logarithmic deadline}

The predicate $\mathcal G$ permits an unrestricted ordinary meeting by
$\lfloor\log_2n\rfloor$. Unlike Gao's predicate it imposes neither
balanced counts nor a before-one condition. The following affine count
pays for unbalanced meetings directly.

\begin{lemma}[Bad original parity residues]\label{gaoexact:bad}
For $0<\delta<1/2$, put
\[
 E_h(\delta)=\{n\ge0:|o_h(n)-h/2|\ge\delta h
            \text{ or }|o_h(n+1)-h/2|\ge\delta h\}.
\]
For integers $N,h\ge1$,
\[
 \frac{\#(E_h(\delta)\cap[0,N))}N
 \le4(1+2^h/N)e^{-2\delta^2h}=:e_h(N,\delta).
\]
\end{lemma}
\begin{proof}
Original parities are conditionally fair by Lemma~\ref{ordinary:parities}.
Hoeffding's bound and a union bound give Haar probability at most
$4e^{-2\delta^2h}$; no independence between the two roots is required.
The first $h$ canonical labels fix both original parity words, so this
is a union of classes modulo $2^h$. The upper bound of
Lemma~\ref{pop:periodic} proves the stated integer estimate.
\end{proof}

\begin{lemma}[Localized unbalanced meetings]\label{gao:unbalanced-count}
For integers $k,h\ge1$, let $L_{k,h}$ be the positive starts with
$\tau_h(n)>k$, $\tau_h(n+1)>k$ and an ordinary meeting at some
$1\le j\le k$ with unequal halving counts. Then $\#L_{k,h}\le k2^h$.
\end{lemma}
\begin{proof}
The first $h$ shortcut words are fixed on each residue modulo $2^h$.
On a class where both completed-grid times exceed $k$, these words fix
both ordinary words through $k$: fewer than $h$ complete blocks have
finished, and a possible final odd branch belongs to a block among the
first $h$. At a fixed $j\le k$, the endpoint maps in the common variable
$n$ have slopes $3^j/6^{H_1}$ and $3^j/6^{H_2}$. Unequal halving counts
make these slopes distinct. Their equality has at most one real solution,
hence at most one positive lift in that residue. Sum over the $2^h$
residues and $k$ times.
\end{proof}

\begin{lemma}[Deadline comparisons]\label{gaoexact:deadline}
\begin{enumerate}
\item If $n>2^{h-1}$, $A(n)\le h$, $n\notin E_h(\delta)$, and
$(3/2+\delta)h\le\lfloor\log_2n\rfloor$, then $n\in\mathcal G$.
Its first balanced witness also precedes either root's visit to one.
\item If $n\ge1$, $A(n)>h$, $n\notin E_h(\delta)$, and
$(3/2-\delta)h>k\ge\lfloor\log_2n\rfloor$, no balanced merger
occurs by $k$. Any witness to $n\in\mathcal G$ belongs to $L_{k,h}$.
\end{enumerate}
\end{lemma}
\begin{proof}
In the first case, Lemma~\ref{ordinary:exactclock} gives
$B(n)\le\tau_h(n)-1<(3/2+\delta)h$. Each prefix value on either root
is at least $n/2^{A(n)-1}>1$, since each path has at most $A(n)-1$
halvings through $B(n)$.
In the second case the same dictionary gives
$B(n)\ge\tau_h(n)\ge(3/2-\delta)h>k$, interpreting $B=\infty$
when $A=\infty$. Good parities give $\tau_h(n+1)>k$ as well, so any
remaining deadline witness is one of the counted unbalanced meetings.
\end{proof}

\begin{theorem}[Exact logarithmic deficiency]\label{gaoexact:loglimit}
For integer cutoffs $N\to\infty$,
\[
 \sqrt{\log_2N}\,f_N\longrightarrow\sqrt{3/2}\,\kappa.
\]
\end{theorem}
\begin{proof}
Write $L=\log_2N$, $k=\lfloor L\rfloor$. Fix $0<\varepsilon<1/6$ and set
\[
 a_\pm=2/3\pm\varepsilon,\quad h_\pm=\lfloor a_\pm L\rfloor,
 \quad\delta=\varepsilon/4,\quad R_N=\lceil N/L^2\rceil.
\]
One has $(3/2+\delta)a_-<1<(3/2-\delta)a_+$. For sufficiently large
$N$ the linear margins dominate rounding and $2\log_2L$, so
\[
 (3/2-\delta)h_+>k,\qquad
 (3/2+\delta)h_-\le\lfloor L-2\log_2L\rfloor.
\]
For $n\ge R_N$ the latter bound is at most $\lfloor\log_2n\rfloor$,
and $R_N>2^{h_--1}$. Apply Lemma~\ref{gaoexact:deadline}, the absorption
discrepancy of Proposition~\ref{pop:discrepancy}, the bad-word estimate,
and the $k2^{h_+}$ unbalanced-witness count. Discard $[0,R_N)$ for the
upper bound. Zero belongs to the complement for the lower bound. This gives
\[
 S_{h_+}(1)-\frac{(k+1)2^{h_+}}N-e_{h_+}(N,\delta)
 \le f_N\le
 S_{h_-}(1)+\frac{2^{h_-}}N+e_{h_-}(N,\delta)+\frac{R_N}N.
\]
For fixed $\varepsilon$, each error is $o(L^{-1/2})$: the cylinder errors
are $O(L2^{-(1-a_+)L})$, the parity errors decay exponentially in $L$,
and $R_N/N\le L^{-2}+N^{-1}$. The shortcut tail therefore gives
\[
 \frac{\kappa}{\sqrt{2/3+\varepsilon}}
 \le\liminf\sqrt L\,f_N\le\limsup\sqrt L\,f_N
 \le\frac{\kappa}{\sqrt{2/3-\varepsilon}}.
\]
Let $\varepsilon\downarrow0$. This proof uses the unequal-slope count,
with no pointwise deadline bridge as a premise.
\end{proof}

\begin{remark}[The before-one endpoint convention]
Gao excludes one in the ``first $k$ iterations'' \cite[p.~264]{Gao1993}.
If the forbidden prefix omits the merger endpoint, the existence predicate
is unchanged: a balanced merger at one has previous endpoints two on both
paths, since two is the unique positive preimage of one. Deleting that
last halving preserves both endpoint equality and balanced counts and
gives an earlier merger strictly above one. The finite adapter error is
also valid if forbidden visits are read over the entire horizon $k$.
\end{remark}

% ---- exact/descent ----
\section{Early merger and published descent witnesses}\label{sec:windows}

The exact laws also provide short synchronized applications of Korec's
and Tao's published singleton theorems. We first obtain the merger set,
then transfer either descent witness by one deterministic ordering argument.
Natural density one means complement counting function $o(N)$;
logarithmic density one means complementary harmonic sum $o(\log N)$.

\begin{lemma}[Finite block counts]\label{pop:blockcounts}\label{lem:window-count}
Let $I$ consist of $N$ consecutive positive anchors and let $h\ge1$.
For an integer $H\ge0$, the number of anchors whose block
$n,\ldots,n+H$ has an adjacent edge unabsorbed at height $h$ is at most
\[
 H\bigl(NS_h(1)+2^hS_h(1)(1-S_h(1))\bigr).
\]
Every other block agrees at ordinary time $2h$ and all later times.
For $H\ge1$, put
\[
 r_{H,h}(n)=\frac1H\#\{0\le d<H:C^{2h}(n+d)\ne C^{2h}(n)\}.
\]
For an absolute constant $C_0$ supplied by Theorem~\ref{pair:uniform-tail},
\[
 \frac1N\sum_{n\in I}r_{H,h}(n)
 \le C_0\sqrt{\frac{1+\log H}{h}}+\frac{2^h}{N}.
\]
\end{lemma}
\begin{proof}
For each adjacent edge, its translated anchors form an interval of exactly
$N$ integers. Proposition~\ref{pop:discrepancy} bounds the nonabsorbed
count on that interval. Sum over the $H$ edges. Absorbed pairs agree by
ordinary time $2h$ by Lemma~\ref{ordinary:exactclock}, hence still agree
at $2h$; equality along the edges synchronizes the whole block.

For each $1\le d<H$, disagreement at time $2h$ implies nonabsorption
at height $h$. The offset version of the same discrepancy bounds its
count by $NS_h(d)+2^h$. Theorem~\ref{pair:uniform-tail} gives
$S_h(d)\le C_0\sqrt{(1+\log(1+d))/h}$, and $1+d\le H$.
The term $d=0$ is zero. Average these bounds over $d$ and divide by $N$.
No independence of edges or offsets is used.
\end{proof}

\begin{lemma}[Two early-merger regimes]\label{lem:early}
Fix $0<\delta<1$ and set $M_j=2\lfloor\delta j\rfloor$.
\begin{enumerate}
\item For every prescribed integer-valued $L(n)\ge0$ with
$L(n)=o(\sqrt{\log n})$, a set of natural density one has
\[
 C^{M_j}(n+d)=C^{M_j}(n)\quad(0\le d\le L(n)),
 \qquad 2^j\le n<2^{j+1}.
\]
\item For prescribed integers $H_j\ge1$ with $\log H_j=o(j)$, there
are $\varepsilon_j\to0$ and a set of natural density one on which
\[
 \#\{0\le d<H_j:C^{M_j}(n+d)=C^{M_j}(n)\}
 \ge(1-\varepsilon_j)H_j.
\]
\end{enumerate}
The assertions hold for all sufficiently large anchors in their sets.
\end{lemma}
\begin{proof}
Write $X=2^j$, $h=\lfloor\delta j\rfloor$. For the first assertion let
$H_j=\max_{X\le n<2X}L(n)$. The hypothesis gives $H_j=o(\sqrt j)$,
without a monotonicity assumption on $L$. Lemma~\ref{pop:blockcounts}
and $S_h(1)=O(h^{-1/2})$ show that the defective fraction is at most
\[
 H_j\bigl(S_h(1)+2^{h-j}\bigr)=o(1).
\]
For the second assertion the same lemma bounds the mean omitted fraction by
\[
 a_j=C_0\sqrt{\frac{1+\log H_j}{\lfloor\delta j\rfloor}}
             +2^{\lfloor\delta j\rfloor-j}\longrightarrow0.
\]
Markov's inequality with $\varepsilon_j=\sqrt{a_j}$ gives defective
fraction at most $\varepsilon_j$; assign the finitely many initial
$\varepsilon_j$ arbitrarily.

In either case the defective fraction tends to zero in full dyadic shells.
Their union has natural density zero: if each shell beyond a fixed index
has at most $\eta2^j$ defects, summing through the shell containing an
arbitrary cutoff $N$ bounds its counting function by a finite prefix plus
$2\eta N$. Let $\eta\downarrow0$. Equality persists after $M_j$.
\end{proof}

\begin{lemma}[Early-merger witness transfer]\label{descent:transfer}
Suppose a natural-density-one set $E$ has, for each sufficiently large
$n\in E$, a nonnegative integer merger time $m(n)$ and a finite set of offsets
$D(n)\subseteq\N$ containing zero such that
\[
 C^{m(n)}(n+d)=C^{m(n)}(n)\quad(d\in D(n)).
\]
Let $v(n)>0$ satisfy $v(n)\le n2^{-m(n)}$ eventually. Suppose a set $W$
of natural or logarithmic density one supplies an ordinary witness
$t(n)$ with $C^{t(n)}(n)<v(n)$. Then a set of that same density has
\[
 t(n)>m(n),\qquad C^{t(n)}(n+d)=C^{t(n)}(n)<v(n)\quad(d\in D(n)).
\]
Any upper deadline for the singleton witness remains valid.
\end{lemma}
\begin{proof}
For every positive start, $C^i(n)\ge n2^{-i}$ by induction. Thus for
$i\le m(n)$ its value is at least $n2^{-m(n)}\ge v(n)$, forcing
$t(n)>m(n)$. Equality at merger persists until that witness time.
The intersection $E\cap W$ has natural density one when $W$ does.
If $W$ has logarithmic density one, the same holds for the intersection:
for a set with counting function $A(x)=o(x)$, partial summation gives
\[
 \sum_{\substack{n\le N\\n\in A}}\frac1n
 =\frac{A(N)}N+\int_1^N\frac{A(x)}{x^2}\,dx=o(\log N).
\]
So the natural-density-zero complement of $E$ has logarithmic density zero.
\end{proof}

\begin{corollary}[Synchronized Korec descent]\label{thm:power-window}\label{cor:majority-power}
Fix $\log_4 3<c<1$. For every prescribed integer-valued $L(n)\ge0$
with $L(n)=o(\sqrt{\log n})$,
natural-density-almost every positive $n$ admits a positive value $y<n^c$
and one ordinary time $t\le2\log_2n$ with
\[
 C^t(n+d)=y\qquad(0\le d\le L(n)).
\]
For every shellwise integer schedule $H_j\ge1$ with $\log H_j=o(j)$, the same
conclusion holds for at least $(1-\varepsilon_j)H_j$ members of
$[n,n+H_j)$, including its anchor, where $\varepsilon_j\to0$ and the
set of anchors has natural density one.
\end{corollary}
\begin{proof}
Korec's published theorem supplies natural-density-almost every $n$ with
$T^r(n)<n^c$ for some integer $0\le r\le\log_2n$
\cite[Theorem~1 and pp.~86--88]{Korec1994}. The deadline is retained from
his construction: successful inputs satisfy $r=m$ and $m2^m<n$.
Expanding the shortcut word gives ordinary time $t=r+o_r(n)\le2r$.

Choose $0<\delta<(1-c)/2$ and use the relevant regime of
Lemma~\ref{lem:early}, with $m(n)=M_j$. Since $M_j\le2\delta\log_2n$,
\[
 n2^{-M_j}\ge n^{1-2\delta}>n^c\qquad(n>1).
\]
Apply Lemma~\ref{descent:transfer} with $v(n)=n^c$ and with $D(n)$ either
the whole prescribed block or the offsets already equal to its anchor.
This preserves Korec's density and the displayed ordinary deadline.
\end{proof}

\begin{corollary}[Synchronized Tao descent]\label{thm:window}\label{descent:majority-tao}
For every $f(n)\to+\infty$ and every prescribed integer-valued $L(n)\ge0$
with $L(n)=o(\sqrt{\log n})$, logarithmic-density-almost every positive $n$
admits one ordinary time $t$ and positive value $y<f(n)$ with
\[
 C^t(n+d)=y\qquad(0\le d\le L(n)).
\]
For every shellwise integer $H_j\ge1$ with $\log H_j=o(j)$, at least
$(1-\varepsilon_j)H_j$ members of $[n,n+H_j)$ share such a $y$ at one
time, including the anchor, on a set of anchors of logarithmic density one.
In both cases the witness follows the merger deadline
$2\lfloor j/4\rfloor$, where $2^j\le n<2^{j+1}$.
\end{corollary}
\begin{proof}
Tao's published Theorem~1.3 gives, for every real-valued $u(n)\to+\infty$,
a logarithmic-density-one set of positive starts with an ordinary iterate
strictly below $u(n)$ \cite{Tao2022}. Use
$u(n)=\min\{f(n),n^{1/4}\}$ and Lemma~\ref{lem:early} with $\delta=1/4$.
Its merger deadline satisfies $M_j\le(\log_2n)/2$, so
$n2^{-M_j}\ge\sqrt n\ge u(n)$ eventually. Apply
Lemma~\ref{descent:transfer} with $v=u$. Tao supplies no upper deadline
for this small-value witness, and none is asserted here.
\end{proof}

The every-member blocks use adjacent edges; the larger blocks average
fixed offsets and omit a vanishing fraction. For example,
$H_j=\lceil\exp(\sqrt j)\rceil$ permits $\varepsilon_j=O_\delta(j^{-1/8})$.
The singleton imports supply the descent thresholds and densities;
the pair laws supply synchronization. These applications retain exceptional
anchors and, in the larger regime, omitted members.

% ---- exact/closing ----
\section{Scope and open questions}\label{sec:scope}
The tail laws concern Haar or fresh-label laws; the population conclusions
use exact integer cylinders. They give no merger or termination theorem
for every positive integer. A pair may share a nontrivial cycle or an
unbounded future. The growing-block descent applications preserve the
density and witness deadline of their published singleton inputs.

The proof uses fixed observer precision and gives no effective remainder
or numerical value of $\E M$. The harmonic capacity concerns fixed states
and fixed histories. The triadic obstruction is pointwise, whereas its
absorption sufficiency uses fair labels.

\subsection*{Assurance of the results and their inputs}
All proofs in this paper are written in full and have been independently
reviewed. In addition, the first limit of Theorem~\ref{thm:exact},
$\sqrt t\,\PP(A>t)\to2\E M/\sqrt\pi$ with $1\le\E M<\infty$, is
kernel-checked in the author's Lean~4 development, which is not yet
public, in the following form. The declaration
\nolinkurl{Collatz.N04CanonicalHaarTail.canonical_haar_absorption_tail}
states that $\sqrt t$ times the Haar measure of the set of
$x\in\Ztwo$ whose canonical state started from $(0,1)$ is not absorbed
at time $\lfloor t\rfloor$ tends to $2c/\sqrt\pi$, where $c$ is the mean
of the canonical departure count under fair labels;
\nolinkurl{Collatz.N04HaarDepartureCount.canonical_haar_tail_with_haar_departure_mean}
shows that $c$ is the Haar mean of the number $M$ of level-zero
departures along the actual two-adic trajectory, and the bounds
$1\le\E M<\infty$ are checked in the same development. This is a
statement about Haar measure on $\Ztwo$, not about integer counts. The
equality of the canonical absorption time with the balanced merger time
$A$ is Lemma~\ref{pair:dictionary}, a written proof. The remaining limits of
Theorem~\ref{thm:exact} and all other results are written proofs only.
The imported inputs are published: the Terras--Lagarias parity
correspondence \cite{Terras1976,Lagarias1985}, the reflection principle
and central binomial estimates \cite{Feller}, Korec's Theorem~1
\cite{Korec1994} and Tao's Theorem~1.3 \cite{Tao2022}, each used only as
stated in the cited source. Duran's preprint \cite{Duran2026} and the
OEIS entry \cite{OEIS398792} are compared with our results and are not
used in any proof.

Gao's separate question remains open: if adjacent positive starts first
reach one at the same finite ordinary time, must their halving counts up
to that time agree? The exact synchronization rates do not settle it.
Quantitative occupation remainders and computable tail constants are
further questions suggested by the method.

\begin{thebibliography}{99}
\bibitem{Feller}
W.~Feller,
\emph{An Introduction to Probability Theory and Its Applications},
vol.~I, third edition, John Wiley \& Sons, New York, 1968.

\bibitem{Gao1993}
G.-G.~Gao,
\emph{On consecutive numbers of the same height in the Collatz problem},
Discrete Math. \textbf{112} (1993), 261--267.
\href{https://doi.org/10.1016/0012-365X(93)90240-T}
{doi:10.1016/0012-365X(93)90240-T}.

\bibitem{Duran2026}
O.~J.~Said Duran,
\emph{Consecutive integers coalesce with density one under the Collatz map:
a formally verified proof of a conjecture of Gao},
Zenodo preprint, version~1, September~27, 2026.
\href{https://doi.org/10.5281/zenodo.23003524}{doi:10.5281/zenodo.23003524}.

\bibitem{Lagarias1985}
J.~C.~Lagarias,
\emph{The $3x+1$ problem and its generalizations},
Amer. Math. Monthly \textbf{92} (1985), no.~1, 3--23.
\href{https://doi.org/10.2307/2322189}{doi:10.2307/2322189}.

\bibitem{Terras1976}
R.~Terras,
\emph{A stopping time problem on the positive integers},
Acta Arith. \textbf{30} (1976), 241--252.
\href{https://doi.org/10.4064/aa-30-3-241-252}{doi:10.4064/aa-30-3-241-252}.

\bibitem{OEIS398792}
OEIS Foundation Inc.,
\emph{The On-Line Encyclopedia of Integer Sequences}, entry A398792,
\url{https://oeis.org/A398792}, accessed October~8, 2026.

\bibitem{Korec1994}
I.~Korec, \emph{A density estimate for the $3x+1$ problem},
Math. Slovaca \textbf{44} (1994), no.~1, 85--89,
\href{https://dml.cz/bitstream/handle/10338.dmlcz/133225/MathSlov_44-1994-1_8.pdf}{primary article}.

\bibitem{Tao2022}
T.~Tao, \emph{Almost all orbits of the Collatz map attain almost bounded values},
Forum Math. Pi \textbf{10} (2022), e12,
\href{https://doi.org/10.1017/fmp.2022.8}{doi:10.1017/fmp.2022.8}.

\bibitem{GaoDensityOne}
D.~Leen, \emph{Gao's dyadic density-one coalescence conjecture},
preprint, Zenodo, 2026.
\url{https://doi.org/10.5281/zenodo.23268762}.

\bibitem{GaoTao}
D.~Leen, \emph{Quantitative synchronization of Collatz trajectories},
preprint, Zenodo, 2026.
\url{https://doi.org/10.5281/zenodo.23268755}.

\bibitem{GeometryCompanion}
D.~Leen, \emph{Late synchronization and the geometry of Collatz families},
preprint, Zenodo, 2026.
\url{https://doi.org/10.5281/zenodo.23268753}.

\bibitem{Umbrella}
D.~Leen, \emph{The stochastic $3x+1$ model is a theorem: Brownian
structure of Collatz orbits}, preprint, Zenodo, 2026.
\url{https://doi.org/10.5281/zenodo.23268765}.

\end{thebibliography}
\end{document}
